% archon:covers Gluck/Sphere.lean
\chapter{The spherical converse (S², positive curvature)}
\label{chap:sphere}

This chapter formalizes the \emph{spherical} converse to the four vertex theorem in
the positive-curvature (Gluck-style) regime: a prescribed cyclic geodesic-curvature
function on the round sphere $S^2$, with the four-vertex property and a uniform lower
bound forcing confinement to an open hemisphere, is the geodesic curvature of a simple
closed curve contained in that hemisphere. The Euclidean development
(\cref{chap:reduction} and its dependencies, culminating in \texttt{gluck\_converse} and
\texttt{dahlbergConverse}) is complete and is reused, not disturbed. The positive-stage
capstone is named \texttt{sphericalConverse\_pos}; the name \texttt{sphericalConverse} is
reserved for the mandatory mixed-sign stage 2.

All spherical geometry is transported to the \emph{stereographic disk model}: the open
unit disk $\{|z|<1\}\subset\mathbb C$ with the round metric
$g_{S^2}=\dfrac{4}{(1+|z|^2)^2}\,|dz|^2$. Stereographic projection from the missed
antipode is a conformal diffeomorphism onto $S^2$ minus a point, so a curve confined to
$\{|z|<1\}$ is simple on $S^2$ iff its disk representative is simple in $\mathbb C$.

% SOURCE: references/spaceform_notes.md, "THE CRUX AND ITS RESOLUTION" +
%   "H² vs S²" sections (read from references/spaceform_notes.md)
% SOURCE QUOTE: "The tangent-angle gauge reduces closure to an honest 2-D problem.
%   Parametrize by the EUCLIDEAN tangent angle θ in the disk model, z_θ = q(θ)(cos θ,
%   sin θ), q>0. Then κ_E = 1/q and the conformal law solves ALGEBRAICALLY for the
%   speed:  S²:  q(θ) = (1+|z|²) / [2(κ_S(θ) + ⟨z,n⟩)]  (den. >0 ⟺ κ_S > −⟨z,n⟩).
%   Reconstruction becomes a nonlinear 2-D ODE z_θ = q(θ)(cos θ, sin θ) on the disk.
%   The tangent angle increases once around ⇒ TURNING NUMBER +1 IS BUILT IN ⇒ the
%   rotational/framing dimension is discharged by construction ..., leaving only the
%   2-D point-closure error E(z₀)=z(2π;z₀)−z₀ ∈ ℝ². Planar winding/degree then applies."

\section{The conformal geodesic-curvature law}
\label{sec:sphere_conformal}

Write $z=(x,y)\in\mathbb R^2\cong\mathbb C$ for a point of the disk and let a curve have
unit Euclidean tangent $T=e^{i\varphi}$ at inclination angle $\varphi$. We use the
\emph{outward} (right) unit normal
\[
  n=-iT=e^{i(\varphi-\pi/2)}=(\sin\varphi,-\cos\varphi),\qquad
  \langle z,n\rangle=x\sin\varphi-y\cos\varphi .
\]
For the isothermal factor $e^{2\lambda}=4/(1+|z|^2)^2$ one has
$\lambda=\ln 2-\ln(1+|z|^2)$ and $\nabla\lambda=-2z/(1+|z|^2)$, so
$\partial_n\lambda=\langle\nabla\lambda,n\rangle=-2\langle z,n\rangle/(1+|z|^2)$. The
conformal transformation law $\kappa_g=e^{-\lambda}(\kappa_E+\partial_n\lambda)$ gives the
fundamental identity used throughout:
\[
  \kappa_S \;=\; \frac{1+|z|^2}{2}\,\kappa_E \;-\; \langle z,n\rangle .
\]
Here $\kappa_E$ is the Euclidean signed curvature of the disk representative and
$\kappa_S$ its spherical geodesic curvature. In the tangent-angle gauge (\cref{sec:sphere_defs})
$\kappa_E=\varphi'/\|z'\|$, so the law rearranges to the \emph{speed relation}
\[
  \frac{1+|z|^2}{2}\,\varphi' \;=\; \bigl(\kappa_S+\langle z,n\rangle\bigr)\,\|z'\| ,
  \qquad\text{gauge speed}\quad q=\frac{1+|z|^2}{2\bigl(\kappa_S+\langle z,n\rangle\bigr)} .
\]

\paragraph{Geodesic-circle sanity anchor (decisive sign check).}
A counterclockwise Euclidean circle $|z|=r$ has $\kappa_E=1/r$ and outward normal
$n=z/r$, so $\langle z,n\rangle=r$ and
$\kappa_S=\dfrac{1+r^2}{2r}-r=\dfrac{1-r^2}{2r}=\cot\rho$ with $r=\tan(\rho/2)$ — exactly
the geodesic curvature of a spherical circle of intrinsic radius $\rho$ (and $0$ on the
equator $\rho=\pi/2$). The gauge speed is then $q=r$. The opposite sign choice would give
$(1+3r^2)/(2r)\ne\cot\rho$, so \emph{the outward normal with a minus sign is the correct
convention}; any Lean formalization must reproduce $q(\text{circle }r)=r$.

\paragraph{Lean encoding of the normal term.}
In Lean the left normal $iT$ is the natural object: $\langle z,iT\rangle=\langle z,
i\,e^{i\varphi}\rangle=-x\sin\varphi+y\cos\varphi=-\langle z,n\rangle$. Hence the speed
relation is encoded with $\langle z,n\rangle$ replaced by $-\langle z,iT\rangle$:
\[
  \frac{1+\|z\|^2}{2}\,\varphi'
    =\bigl(\kappa-\langle z,\, i\,e^{i\varphi}\rangle_{\mathbb R}\bigr)\,\|z'\| .
\]
The project supplies the rewrite lemma
$\langle z, i\,e^{i\varphi}\rangle_{\mathbb R}=-(\operatorname{Re}z)\sin\varphi
+(\operatorname{Im}z)\cos\varphi$ so both forms are interchangeable.

\section{Definition layer}
\label{sec:sphere_defs}

\begin{lemma}\leanok
[normal-term coordinate identity]
  \label{lem:sphere_normal_inner_eq}
  \lean{Gluck.sphereNormal_inner_eq}
  For $z\in\mathbb C$ and $\varphi\in\mathbb R$,
  $\langle z,\, i\,e^{i\varphi}\rangle_{\mathbb R}
   =-(\operatorname{Re}z)\sin\varphi+(\operatorname{Im}z)\cos\varphi$.
  Equivalently $\langle z, i\,e^{i\varphi}\rangle_{\mathbb R}=-\langle z,n\rangle$ for the
  outward normal $n=-i\,e^{i\varphi}$, so the defining coefficient of
  \cref{def:realizes_spherical_curvature} is
  $\kappa-\langle z,i\,e^{i\varphi}\rangle_{\mathbb R}=\kappa+\langle z,n\rangle$.
\end{lemma}

\begin{proof}

\leanok
  Expand the real inner product $\langle w,v\rangle_{\mathbb R}=(v\,\overline w).\mathrm{re}$
  with $v=i\,e^{i\varphi}$, using
  $(e^{i\varphi}).\mathrm{re}=\cos\varphi$, $(e^{i\varphi}).\mathrm{im}=\sin\varphi$; the
  real and imaginary multiplications collapse to
  $-(\operatorname{Re}z)\sin\varphi+(\operatorname{Im}z)\cos\varphi$.
\end{proof}

\begin{definition}\leanok
[realizes spherical curvature]
  \label{def:realizes_spherical_curvature}
  \lean{Gluck.RealizesSphericalCurvature}
  A curve $z:\mathbb R\to\mathbb C$ \emph{realizes the spherical curvature function}
  $\kappa:\mathbb R\to\mathbb R$ when it is $C^1$, regular ($z'(t)\neq 0$ for all $t$),
  confined to the open disk ($\|z(t)\|<1$ for all $t$), and there is a differentiable
  tangent-angle function $\varphi:\mathbb R\to\mathbb R$ with, for all $t$,
  \[
    z'(t)=\|z'(t)\|\,e^{i\varphi(t)},\qquad
    \frac{1+\|z(t)\|^2}{2}\,\varphi'(t)
      =\bigl(\kappa(t)-\langle z(t),\,i\,e^{i\varphi(t)}\rangle_{\mathbb R}\bigr)\,\|z'(t)\| .
  \]
  By \cref{sec:sphere_conformal} the coefficient $\kappa-\langle z,iT\rangle
  =\kappa+\langle z,n\rangle$ is the reciprocal-speed weight of the conformal law, so this
  says precisely that $\kappa$ is the spherical geodesic curvature of $z$. It mirrors the
  Euclidean \cref{def:realizes_curvature} exactly: it is \emph{speed-free} (only the angle
  $\varphi$, not $z$, must be differentiable; the speed $\|z'\|$ is never normalized to arc
  length), so it is meaningful for a merely continuous $\kappa$ and does not reintroduce
  the three-dimensional spherical frame problem.
\end{definition}

\begin{definition}\leanok
[spherical four-vertex condition, positive stage]
  \label{def:sphere_four_vertex}
  \lean{Gluck.SphereFourVertex}
  \uses{def:four_vertex_condition, def:is_curvature_function}
  A function $\kappa:\mathbb R\to\mathbb R$ satisfies the \emph{positive-stage spherical
  four-vertex condition} when it is a curvature function (\cref{def:is_curvature_function}:
  continuous, $2\pi$-periodic, strictly positive) and satisfies the value-separated Euclidean
  four-vertex condition \cref{def:four_vertex_condition} — exactly the hypothesis package of
  the Euclidean \texttt{gluck\_converse}. No confinement bound is bundled: by compactness a
  uniform bound $R$ with $0<R<1$ and $\kappa>R$ is automatic
  (\cref{lem:sphere_curvature_lower_bound}), and a bundled witness would be gratuitous
  interface drift from the Euclidean theorem without carrying any datum the realization
  layer consumes (audit finding, iter-043). The bound $\kappa>R$ is what makes the
  denominator $\kappa_S+\langle z,n\rangle$ strictly positive on $\{|z|\le R\}$
  (\cref{lem:sphere_denom_pos}); it is the positive-stage stand-in for the sign-changing
  admissibility $\kappa_S>-\langle z,n\rangle$ of stage 2.
\end{definition}

\begin{lemma}\leanok
[uniform curvature lower bound]
  \label{lem:sphere_curvature_lower_bound}
  \lean{Gluck.exists_curvature_lower_bound}
  \uses{def:is_curvature_function}
  If $\kappa$ is a curvature function (\cref{def:is_curvature_function}), then there is $R$
  with $0<R$, $R<1$, and $R<\kappa(\theta)$ for all $\theta$.
\end{lemma}

\begin{proof}

\leanok
  $\kappa$ is continuous on the compact interval $[0,2\pi]$, so it attains there a minimum
  value $m=\kappa(\theta_0)>0$. Set $R=\min(m,1)/2$. Then $0<R$, $R<1$, and
  $R<m\le\kappa(\theta)$ for every $\theta\in[0,2\pi]$; an arbitrary $\theta$ reduces to
  this window by $2\pi$-periodicity (subtract the integer multiple
  $2\pi\lfloor\theta/(2\pi)\rfloor$).
\end{proof}

\begin{theorem}\leanok
[spherical converse, positive stage]
  \label{thm:spherical_converse_pos}
  \lean{Gluck.sphericalConverse_pos}
  \uses{def:realizes_spherical_curvature, def:sphere_four_vertex,
        lem:spherical_speed_solve, lem:reconstruction_ode,
        lem:spherical_endpoint_winding, lem:realizes_spherical_comp,
        lem:spherical_simplicity, lem:spherical_circle_realizes}
  If $\kappa$ satisfies the positive-stage spherical four-vertex condition
  (\cref{def:sphere_four_vertex}), then there is a simple closed curve $z$ confined to the
  open disk that realizes $\kappa$ as its spherical geodesic curvature:
  \[
    \texttt{SphereFourVertex }\kappa \;\Longrightarrow\;
    \exists\,z,\ \texttt{IsSimpleClosed }z\ \wedge\ \texttt{RealizesSphericalCurvature }z\,\kappa .
  \]
  This is the same conclusion shape as the Euclidean \texttt{gluck\_converse}, with
  \texttt{RealizesCurvature} replaced by its spherical analogue.
\end{theorem}

\begin{proof}
  \uses{lem:reconstruction_ode, lem:spherical_endpoint_winding,
        lem:spherical_simplicity, lem:spherical_circle_realizes,
        lem:realizes_spherical_comp, lem:exists_c1_circle_inverse}
  Split on the four-vertex disjunction (\cref{def:four_vertex_condition}). The
  \emph{constant} branch is discharged directly by the explicit centered circle
  (\cref{lem:spherical_circle_realizes}), with no flow machinery. For the value-separated
  (non-constant) branch: \cref{lem:spherical_endpoint_winding} supplies $R<1$,
  $\delta>0$, a circle reparametrization $h_1$, and a \emph{closed, admissible}
  trajectory $z$ of the $(\kappa\circ h_1,R,\delta)$-truncated flow.
  \cref{lem:reconstruction_ode} (with curvature function $\kappa\circ h_1$) upgrades its
  periodic extension to a closed $C^1$ curve confined to $\{|z|\le R\}\subset\{|z|<1\}$
  realizing $\kappa\circ h_1$ in the gauge $\varphi(\theta)=\theta$.
  \cref{lem:realizes_spherical_comp} with $\psi=h_1^{-1}$ pulls this back to a curve
  realizing $\kappa$ (closedness and injectivity-on-a-period transfer along the
  circle-equivariant $h_1^{-1}$, as in the Euclidean
  \texttt{isSimpleClosed\_comp}). \cref{lem:spherical_simplicity} gives simplicity in
  the disk, hence on $S^2$. This mirrors how \texttt{dahlbergConverse} assembles its
  closing parameter, simplicity transport, and reconstruction. \emph{(Assembled formally
  in S2-E once the S2-D layer exists in Lean.)}
\end{proof}

\section{Proof arc}
\label{sec:sphere_arc}

The proof follows the Euclidean Route A architecture, with the reconstruction integral
replaced by a reconstruction ODE. The architecture was pinned at iteration 046
(design analysis + external consult, load-bearing computations re-verified in-project):
\begin{enumerate}
\item \textbf{S2-B (truncated flow, unconditional).} The gauge speed is truncated
  algebraically (\cref{def:truncated_speed}) so the ODE field becomes globally
  Lipschitz and bounded; global existence on $[0,2\pi]$, uniqueness, and continuous
  dependence on the initial point are then unconditional, with no confinement
  prerequisite. Only the closed curve found at the end must be shown, a posteriori, to
  be \emph{admissible} (both truncations inactive), whereupon it is an honest spherical
  curve.
\item \textbf{S2-C (admissibility by margin transport).} The confinement mechanism is
  \emph{perturbative}: an explicit model trajectory (a spherical bicircle made of
  circular arcs, \cref{lem:constant_curvature_arc}) is admissible with quantitative
  margins, and a Gr\"onwall estimate with $L^1$-in-$\theta$ drive
  (\cref{lem:invariant_admissible_domain}) transports the margins to every trajectory
  whose curvature is $L^1$-close and whose start is near the model start. The
  Dahlberg-style reparametrization already formalized in \cref{chap:dahlberg_step1}
  supplies arbitrarily small $L^1$ distance to a two-value step model, so no
  uniform-in-$\kappa_0$ direct confinement estimate is needed (none exists: no disk is
  flow-invariant, see the discussion in \cref{lem:invariant_admissible_domain}).
\item \textbf{S2-D (endpoint map and winding).} The degree argument runs over a compact
  $2$-disk of \emph{initial points} $z_0$, but its computable model is the symmetric
  two-value ($\pi$-periodic) step curvature, exploited through the half-turn
  equivariance $z\mapsto -z$, $\theta\mapsto\theta+\pi$
  (\cref{lem:flow_half_turn_equivariance}). The constant-curvature model is
  \emph{unusable}: for constant $\kappa$ every admissible trajectory closes (all
  spherical circles), so the endpoint error vanishes identically and carries no degree
  — see the DO-NOT recorded in \cref{lem:spherical_endpoint_winding}.
\end{enumerate}
The speed layer below is formalized; the truncated-flow layer is the current prover
objective; S2-C/S2-D blocks state the pinned design with honest flags on the two
remaining open sub-questions (both in \cref{lem:spherical_endpoint_winding}).

\begin{definition}\leanok
[gauge speed]
  \label{def:spherical_speed}
  \lean{Gluck.sphericalSpeed}
  \uses{lem:sphere_normal_inner_eq}
  For $\kappa:\mathbb R\to\mathbb R$, $\theta\in\mathbb R$ and $z\in\mathbb C$ set
  \[
    q_\kappa(\theta,z)\;=\;\frac{1+\|z\|^2}
      {2\bigl(\kappa(\theta)-\langle z,\,i\,e^{i\theta}\rangle_{\mathbb R}\bigr)} .
  \]
  By \cref{lem:sphere_normal_inner_eq} the bracket equals
  $\kappa(\theta)+\langle z,n\rangle$ for the outward normal $n=-i\,e^{i\theta}$ of the
  gauge $\varphi(\theta)=\theta$, so $q_\kappa$ is the algebraic solution of the speed
  relation of \cref{def:realizes_spherical_curvature} for the speed $\|z'\|$. It is a total
  function of the junk-value kind: where the bracket vanishes it returns
  division-by-zero junk, and every lemma about it carries the admissibility hypothesis
  $\|z\|\le R<\kappa(\theta)$.
\end{definition}

\begin{lemma}\leanok
[denominator positivity]
  \label{lem:sphere_denom_pos}
  \lean{Gluck.sphere_denom_pos}
  \uses{lem:sphere_normal_inner_eq}
  If $\|z\|\le R$ and $R<\kappa(\theta)$, then
  $\kappa(\theta)-\langle z,\,i\,e^{i\theta}\rangle_{\mathbb R}\;\ge\;\kappa(\theta)-R\;>\;0$.
\end{lemma}

\begin{proof}

\leanok
  Cauchy--Schwarz gives
  $\bigl|\langle z,\,i\,e^{i\theta}\rangle_{\mathbb R}\bigr|\le\|z\|\cdot\|i\,e^{i\theta}\|
  =\|z\|\le R$, since $\|i\,e^{i\theta}\|=|i|\cdot|e^{i\theta}|=1$. Hence
  $\kappa(\theta)-\langle z, i\,e^{i\theta}\rangle_{\mathbb R}\ge\kappa(\theta)-R>0$.
\end{proof}

\begin{lemma}\leanok
[positive speed / algebraic solve]
  \label{lem:spherical_speed_solve}
  \lean{Gluck.sphericalSpeed_pos}
  \uses{def:spherical_speed, lem:sphere_denom_pos, lem:sphere_curvature_lower_bound}
  If $\|z\|\le R$ and $R<\kappa(\theta)$, then $q_\kappa(\theta,z)>0$ (the hypotheses are
  supplied under \cref{def:sphere_four_vertex} by \cref{lem:sphere_curvature_lower_bound}).
  Positivity of $q_\kappa$ makes the tangent turn once around as $\theta$ runs over
  $[0,2\pi]$, so along any single closed reconstructed curve the turning number is $+1$ by
  construction and the rotational/framing dimension is discharged (the exact role
  $\oint\kappa=2\pi$ plays in the plane). This fixes the framing dimension only; it does not
  by itself give point closure, confinement, or the boundary winding of the endpoint map.
\end{lemma}

\begin{proof}

\leanok
  The numerator $1+\|z\|^2$ is strictly positive, and the denominator
  $2(\kappa(\theta)-\langle z, i\,e^{i\theta}\rangle_{\mathbb R})$ is strictly positive by
  \cref{lem:sphere_denom_pos}; a quotient of positives is positive.
\end{proof}

\begin{lemma}\leanok
[speed continuity]
  \label{lem:spherical_speed_continuous}
  \lean{Gluck.sphericalSpeed_continuousOn}
  \uses{def:spherical_speed, lem:sphere_denom_pos}
  If $\kappa$ is continuous and $R<\kappa(\theta)$ for all $\theta$, then
  $(\theta,z)\mapsto q_\kappa(\theta,z)$ is continuous on the slab
  $\{(\theta,z):\|z\|\le R\}\subseteq\mathbb R\times\mathbb C$.
\end{lemma}

\begin{proof}

\leanok
  The numerator $(\theta,z)\mapsto 1+\|z\|^2$ is continuous, and the denominator
  $(\theta,z)\mapsto 2(\kappa(\theta)-\langle z, i\,e^{i\theta}\rangle_{\mathbb R})$ is
  continuous ($\kappa$ is continuous; the inner product is continuous bilinear and
  $\theta\mapsto i\,e^{i\theta}$ is continuous). On the slab the denominator is nonvanishing
  by \cref{lem:sphere_denom_pos}, so the quotient is continuous there.
\end{proof}

\begin{lemma}\leanok
[speed periodicity]
  \label{lem:spherical_speed_periodic}
  \lean{Gluck.sphericalSpeed_periodic}
  \uses{def:spherical_speed}
  If $\kappa$ is $2\pi$-periodic, then $q_\kappa(\theta+2\pi,z)=q_\kappa(\theta,z)$ for all
  $\theta$ and $z$.
\end{lemma}

\begin{proof}

\leanok
  $\kappa(\theta+2\pi)=\kappa(\theta)$ by hypothesis, and
  $e^{i(\theta+2\pi)}=e^{i\theta}e^{2\pi i}=e^{i\theta}$ since $e^{2\pi i}=1$; both
  numerator and denominator of $q_\kappa$ are unchanged.
\end{proof}

\begin{lemma}

\leanok
[geodesic-circle sanity anchor]
  \label{lem:spherical_speed_circle}
  \lean{Gluck.sphericalSpeed_circle}
  \uses{def:spherical_speed}
  For $r>0$ let $\kappa\equiv(1-r^2)/(2r)$ be the constant geodesic curvature of the
  spherical circle of stereographic radius $r$, and let $z=-r\,i\,e^{i\theta}$ be the point
  of that circle at tangent angle $\theta$ (counterclockwise gauge). Then
  $q_\kappa(\theta,z)=r$: the gauge speed reproduces the Euclidean radius. This is the
  decisive machine-checked \emph{sign certificate} of \cref{sec:sphere_conformal} — with
  the opposite sign convention the value would be $\ne r$ and the build would fail.
\end{lemma}

\begin{proof}

\leanok
  $\|z\|=r$ and $\langle z,\,i e^{i\theta}\rangle=-r\,\|i e^{i\theta}\|^2=-r$, so the
  denominator is $2(\kappa+r)=2\bigl(\tfrac{1-r^2}{2r}+r\bigr)=\tfrac{1+r^2}{r}$ and
  $q=\dfrac{1+r^2}{(1+r^2)/r}=r$.
\end{proof}

\section{Truncated flow layer (S2-B)}
\label{sec:sphere_flow}

\emph{Design decision (iter-046).} The reconstruction field is truncated
\emph{algebraically} — the norm clamped in the numerator, the denominator clamped from
below — so that it becomes globally defined, bounded, and globally Lipschitz in $z$. All
flow machinery (existence on $[0,2\pi]$, uniqueness, continuous dependence, the endpoint
map) is then \emph{unconditional}: no confinement lemma is needed to run the degree
argument. Admissibility is re-imposed a posteriori, only on the single closed trajectory
the winding argument produces (\cref{sec:sphere_admissible}). This removes the
continuation/stitching bookkeeping entirely: one Picard--Lindel\"of application on a
large ball covers $[0,2\pi]$.

\begin{definition}

\leanok
[truncated gauge speed]
  \label{def:truncated_speed}
  \lean{Gluck.truncatedSpeed}
  \uses{def:spherical_speed}
  For parameters $R,\delta$ set
  \[
    \hat q_{\kappa,R,\delta}(\theta,z)\;=\;
      \frac{1+\min(\|z\|,R)^2}
        {2\,\max\bigl(\kappa(\theta)-\langle z,\,i\,e^{i\theta}\rangle_{\mathbb R},\ \delta\bigr)}.
  \]
  On the \emph{admissible set}
  $\{(\theta,z):\|z\|\le R\ \wedge\ \kappa(\theta)-\langle z,i e^{i\theta}\rangle\ge\delta\}$
  both clamps are inactive and $\hat q=q_\kappa$; off it, $\hat q$ is a globally tame
  surrogate.
\end{definition}

\begin{lemma}

\leanok
[truncated speed agrees on the admissible set]
  \label{lem:truncated_speed_eq}
  \lean{Gluck.truncatedSpeed_eq}
  \uses{def:truncated_speed, def:spherical_speed}
  If $\|z\|\le R$ and $\kappa(\theta)-\langle z,i e^{i\theta}\rangle_{\mathbb R}\ge\delta$,
  then $\hat q_{\kappa,R,\delta}(\theta,z)=q_\kappa(\theta,z)$.
\end{lemma}

\begin{proof}

\leanok
  $\min(\|z\|,R)=\|z\|$ by the first hypothesis and
  $\max(\kappa(\theta)-\langle z,i e^{i\theta}\rangle,\delta)
   =\kappa(\theta)-\langle z,i e^{i\theta}\rangle$ by the second; the two formulas
  coincide.
\end{proof}

\begin{lemma}

\leanok
[truncated speed is positive]
  \label{lem:truncated_speed_pos}
  \lean{Gluck.truncatedSpeed_pos}
  \uses{def:truncated_speed}
  If $\delta>0$ then $\hat q_{\kappa,R,\delta}(\theta,z)>0$ for all $\theta,z$.
\end{lemma}

\begin{proof}

\leanok
  The numerator is $\ge 1>0$ (a square plus one) and the denominator is
  $\ge 2\delta>0$; a quotient of positives is positive.
\end{proof}

\begin{lemma}

\leanok
[truncated speed is bounded]
  \label{lem:truncated_speed_le}
  \lean{Gluck.truncatedSpeed_le}
  \uses{def:truncated_speed}
  If $0\le R$ and $0<\delta$ then
  $\hat q_{\kappa,R,\delta}(\theta,z)\le B:=\dfrac{1+R^2}{2\delta}$ for all $\theta,z$.
\end{lemma}

\begin{proof}

\leanok
  $\min(\|z\|,R)\le R$ bounds the numerator by $1+R^2$, and the denominator is
  $\ge 2\delta$; divide.
\end{proof}

\begin{lemma}

\leanok
[truncated speed is Lipschitz in $z$, uniformly in $\theta$]
  \label{lem:truncated_speed_lipschitz}
  \lean{Gluck.truncatedSpeed_lipschitz}
  \uses{def:truncated_speed}
  If $0\le R$ and $0<\delta$ there is $L\ge 0$ (explicitly
  $L=\tfrac{R}{\delta}+\tfrac{1+R^2}{2\delta^2}$ works) such that for every $\theta$ the
  map $z\mapsto\hat q_{\kappa,R,\delta}(\theta,z)$ is $L$-Lipschitz on all of $\mathbb C$.
\end{lemma}

\begin{proof}

\leanok
  Write $\hat q=N/D$ with $N(z)=1+\min(\|z\|,R)^2$ and
  $D(z)=2\max(\kappa(\theta)-\langle z,i e^{i\theta}\rangle,\delta)$. The norm is
  $1$-Lipschitz, the clamp $x\mapsto\min(x,R)$ is $1$-Lipschitz and bounded by $R$ (for
  $x\ge0$), so $N$ is $2R$-Lipschitz (difference of squares of clamped values, both
  $\le R$) and $1\le N\le 1+R^2$. The linear functional
  $z\mapsto\langle z,i e^{i\theta}\rangle$ is $1$-Lipschitz (Cauchy--Schwarz against a unit
  vector) and $x\mapsto\max(x,\delta)$ is $1$-Lipschitz, so $D$ is $2$-Lipschitz and
  $D\ge 2\delta$. Then
  \[
    \Bigl|\frac{N(z)}{D(z)}-\frac{N(w)}{D(w)}\Bigr|
    \le\frac{|N(z)-N(w)|}{D(w)}+N(z)\,\frac{|D(z)-D(w)|}{D(z)D(w)}
    \le\Bigl(\frac{2R}{2\delta}+\frac{(1+R^2)\cdot2}{4\delta^2}\Bigr)|z-w| .
  \]
  (Formalization route: reciprocal is $\tfrac1{(2\delta)^2}$-Lipschitz on
  $[2\delta,\infty)$ by the mean value inequality; compose and multiply
  bounded-Lipschitz factors.)
\end{proof}

\begin{lemma}

\leanok
[truncated speed is jointly continuous]
  \label{lem:truncated_speed_continuous}
  \lean{Gluck.truncatedSpeed_continuous}
  \uses{def:truncated_speed}
  If $\kappa$ is continuous and $0<\delta$, then
  $(\theta,z)\mapsto\hat q_{\kappa,R,\delta}(\theta,z)$ is continuous on
  $\mathbb R\times\mathbb C$.
\end{lemma}

\begin{proof}

\leanok
  Numerator and denominator are continuous ($\min$, $\max$, the norm, $\kappa$, the inner
  product against the continuous unit field $\theta\mapsto i e^{i\theta}$), and the
  denominator never vanishes (it is $\ge 2\delta>0$), so the quotient is continuous
  everywhere — no slab restriction, unlike \cref{lem:spherical_speed_continuous}.
\end{proof}

\begin{definition}

\leanok
[truncated reconstruction field]
  \label{def:truncated_field}
  \lean{Gluck.truncatedField}
  \uses{def:truncated_speed}
  $F_{\kappa,R,\delta}(\theta,z)=\hat q_{\kappa,R,\delta}(\theta,z)\,e^{i\theta}
   \in\mathbb C$. Since $\|e^{i\theta}\|=1$, the field inherits the bound
  $\|F\|\le B$ of \cref{lem:truncated_speed_le} and the Lipschitz constant $L$ of
  \cref{lem:truncated_speed_lipschitz} (the difference at fixed $\theta$ is
  $(\hat q(z)-\hat q(w))\,e^{i\theta}$).
\end{definition}

\begin{lemma}\leanok
[field norm equals speed]
  \label{lem:norm_truncated_field}
  \lean{Gluck.norm_truncatedField}
  \uses{def:truncated_field, lem:truncated_speed_pos}
  For $0<\delta$,
  $\|F_{\kappa,R,\delta}(\theta,z)\|=\hat q_{\kappa,R,\delta}(\theta,z)$ for all
  $\theta,z$.
\end{lemma}

\begin{proof}

\leanok
  $\|\hat q\,e^{i\theta}\|=|\hat q|\,\|e^{i\theta}\|=|\hat q|=\hat q$, since the unit
  field has norm one and $\hat q>0$ by \cref{lem:truncated_speed_pos}.
\end{proof}

\begin{lemma}\leanok
[field Lipschitz bound]
  \label{lem:truncated_field_lipschitz}
  \lean{Gluck.truncatedField_lipschitz}
  \uses{def:truncated_field, lem:truncated_speed_lipschitz}
  For $0\le R$, $0<\delta$ there is a constant $L$ (the one of
  \cref{lem:truncated_speed_lipschitz}) such that for every $\theta$ the map
  $z\mapsto F_{\kappa,R,\delta}(\theta,z)$ is $L$-Lipschitz.
\end{lemma}

\begin{proof}

\leanok
  At fixed $\theta$,
  $F(\theta,z)-F(\theta,w)=(\hat q(\theta,z)-\hat q(\theta,w))\,e^{i\theta}$, whose norm
  is $|\hat q(\theta,z)-\hat q(\theta,w)|\le L\|z-w\|$ by
  \cref{lem:truncated_speed_lipschitz}.
\end{proof}

\begin{lemma}\leanok
[field joint continuity]
  \label{lem:truncated_field_continuous}
  \lean{Gluck.truncatedField_continuous}
  \uses{def:truncated_field, lem:truncated_speed_continuous}
  If $\kappa$ is continuous and $0<\delta$, then
  $(\theta,z)\mapsto F_{\kappa,R,\delta}(\theta,z)$ is continuous on
  $\mathbb R\times\mathbb C$.
\end{lemma}

\begin{proof}

\leanok
  Scalar multiplication of the jointly continuous truncated speed
  (\cref{lem:truncated_speed_continuous}) by the continuous unit field
  $\theta\mapsto e^{i\theta}$.
\end{proof}

\begin{lemma}

\leanok
[Picard--Lindel\"of package for the truncated field]
  \label{lem:truncated_field_picard}
  \lean{Gluck.truncatedField_isPicardLindelof}
  \uses{def:truncated_field, lem:truncated_speed_le, lem:truncated_speed_lipschitz,
        lem:truncated_speed_continuous, lem:norm_truncated_field,
        lem:truncated_field_lipschitz, lem:truncated_field_continuous}
  Let $\kappa$ be continuous, $0\le R$, $0<\delta$, and fix an initial-disk radius
  $r_0\ge0$. With $a=r_0+2\pi B+1$, the field $F_{\kappa,R,\delta}$ satisfies the
  Picard--Lindel\"of hypothesis package on the time interval $[0,2\pi]$ with initial time
  $t_0=0$, center $x_0=0$, inner radius $r_0$ and outer radius $a$: it is Lipschitz in the
  space variable with a constant uniform in $\theta$, jointly continuous, norm-bounded by
  $B$, and satisfies the budget condition $B\cdot 2\pi\le a-r_0$.
\end{lemma}

\begin{proof}

\leanok
  All four items are \cref{lem:truncated_speed_lipschitz},
  \cref{lem:truncated_speed_continuous}, \cref{lem:truncated_speed_le}, and the arithmetic
  choice of $a$ respectively. Because the truncated field is bounded and Lipschitz on all
  of $\mathbb C$, restriction to the working ball loses nothing and the budget can always
  be met by enlarging $a$ — this is the payoff of truncation: one application covers
  $[0,2\pi]$ with no continuation argument.
\end{proof}

\begin{lemma}

\leanok
[global flow with continuous dependence]
  \label{lem:spherical_flow_exists}
  \lean{Gluck.exists_sphericalFlow}
  \uses{lem:truncated_field_picard}
  Under the hypotheses of \cref{lem:truncated_field_picard} there is a map
  $\alpha:\mathbb C\times\mathbb R\to\mathbb C$ such that for every $z_0$ in the closed
  disk $\|z_0\|\le r_0$: $\alpha(z_0,0)=z_0$; $\theta\mapsto\alpha(z_0,\theta)$ has on
  $[0,2\pi]$ the derivative $F_{\kappa,R,\delta}(\theta,\alpha(z_0,\theta))$ (within the
  interval); and $\alpha$ is continuous on
  $\{\|z_0\|\le r_0\}\times[0,2\pi]$.
\end{lemma}

\begin{proof}

\leanok
  This is the flow form of the Picard--Lindel\"of theorem
  (in Mathlib:
  \texttt{IsPicardLindelof.exists\_forall\_mem\_closedBall\_eq\_hasDerivWithinAt\_continuousOn}),
  applied to the package of \cref{lem:truncated_field_picard}. It returns a single
  $\alpha$ defined on the product with exactly the three listed properties.
\end{proof}

\begin{definition}

\leanok
[spherical flow]
  \label{def:spherical_flow}
  \lean{Gluck.sphericalFlow}
  \uses{lem:spherical_flow_exists}
  Fix, once per parameter tuple $(\kappa,R,\delta,r_0)$, a choice
  $\Phi=\Phi_{\kappa,R,\delta,r_0}:\mathbb C\times\mathbb R\to\mathbb C$ of the map
  supplied by \cref{lem:spherical_flow_exists}. (The whole flow is chosen at once — not a
  separate choice per initial point — so downstream continuity statements can consume it.)
\end{definition}

\begin{lemma}

\leanok
[flow specification]
  \label{lem:spherical_flow_spec}
  \lean{Gluck.sphericalFlow_spec}
  \uses{def:spherical_flow}
  For $\|z_0\|\le r_0$: $\Phi(z_0,0)=z_0$ and $\theta\mapsto\Phi(z_0,\theta)$ solves
  $z'=F_{\kappa,R,\delta}(\theta,z)$ on $[0,2\pi]$ (derivative within the interval).
\end{lemma}

\begin{proof}

\leanok
  Unfold the choice of \cref{def:spherical_flow}.
\end{proof}

\begin{lemma}

\leanok
[flow continuity]
  \label{lem:spherical_flow_continuousOn}
  \lean{Gluck.sphericalFlow_continuousOn}
  \uses{def:spherical_flow}
  $\Phi$ is continuous on $\{\|z_0\|\le r_0\}\times[0,2\pi]$.
\end{lemma}

\begin{proof}

\leanok
  Unfold the choice of \cref{def:spherical_flow}.
\end{proof}

\begin{lemma}

\leanok
[flow uniqueness]
  \label{lem:spherical_flow_unique}
  \lean{Gluck.sphericalFlow_unique}
  \uses{def:spherical_flow, lem:truncated_speed_lipschitz, lem:spherical_flow_spec}
  If $g:[0,2\pi]\to\mathbb C$ solves $g'=F_{\kappa,R,\delta}(\theta,g(\theta))$ on
  $[0,2\pi]$ with $g(0)=z_0$, $\|z_0\|\le r_0$, then $g=\Phi(z_0,\cdot)$ on $[0,2\pi]$.
\end{lemma}

\begin{proof}

\leanok
  Both solve the same ODE with the same initial value; the field is Lipschitz in the space
  variable uniformly in time (\cref{lem:truncated_speed_lipschitz}), so the standard
  uniqueness theorem for ODEs (Gr\"onwall; in Mathlib
  \texttt{ODE\_solution\_unique\_of\_mem\_Icc}) applies. Uniqueness is what later
  identifies explicitly constructed trajectories — circular arcs, reflected trajectories —
  with the flow.
\end{proof}

\begin{lemma}\leanok
[two-solution uniqueness on a subinterval]
  \label{lem:truncated_field_solution_unique}
  \lean{Gluck.truncatedField_solution_unique}
  \uses{def:truncated_field, lem:truncated_field_lipschitz}
  Let $0\le R$, $0<\delta$ and $0\le T$. If $g_1,g_2:[0,T]\to\mathbb C$ both solve
  $g'=F_{\kappa,R,\delta}(\theta,g(\theta))$ on $[0,T]$ and $g_1(0)=g_2(0)$, then
  $g_1=g_2$ on $[0,T]$. (Unlike \cref{lem:spherical_flow_unique} this compares two
  arbitrary solutions, on an arbitrary compact interval $[0,T]$ — not necessarily
  $[0,2\pi]$ — and makes no reference to the chosen flow $\Phi$; the field is defined
  and Lipschitz for all times, so no time restriction is needed.)
\end{lemma}

\begin{proof}

\leanok
  The field is Lipschitz in the space variable uniformly in time
  (\cref{lem:truncated_field_lipschitz}) and each $g_i$ is continuous on $[0,T]$ (it has
  a derivative there), so the standard two-solution uniqueness theorem applies on
  $[0,T]$ with initial time at the left endpoint — the same Gr\"onwall-type argument as
  \cref{lem:spherical_flow_unique}, with $2\pi$ replaced by $T$.
\end{proof}

\begin{definition}

\leanok
[spherical endpoint map]
  \label{def:spherical_endpoint}
  \lean{Gluck.sphericalEndpoint}
  \uses{def:spherical_flow}
  $E(z_0)=\Phi(z_0,2\pi)-z_0$. A zero of $E$ is a closed trajectory of the truncated
  flow.
\end{definition}

\begin{lemma}

\leanok
[endpoint map continuity]
  \label{lem:spherical_endpoint_continuousOn}
  \lean{Gluck.sphericalEndpoint_continuousOn}
  \uses{def:spherical_endpoint, lem:spherical_flow_continuousOn}
  $E$ is continuous on the closed disk $\{\|z_0\|\le r_0\}$.
\end{lemma}

\begin{proof}

\leanok
  $z_0\mapsto\Phi(z_0,2\pi)$ is continuous by \cref{lem:spherical_flow_continuousOn}
  (restriction of a jointly continuous map to a time slice), and subtracting the identity
  preserves continuity.
\end{proof}

\section{Admissibility and truncation removal (S2-C)}
\label{sec:sphere_admissible}

\begin{lemma}\leanok
[curvature sensitivity of the truncated speed]
  \label{lem:truncated_speed_sub_le}
  \lean{Gluck.truncatedSpeed_sub_le}
  \uses{def:truncated_speed}
  Let $0\le R$, $0<\delta$ and set $M=\tfrac{1+R^2}{2\delta^2}$. Then for all
  $\kappa,\kappa^\ast:\mathbb R\to\mathbb R$, $\theta$ and $z$,
  \[
    \bigl|\hat q_{\kappa,R,\delta}(\theta,z)-\hat q_{\kappa^\ast,R,\delta}(\theta,z)\bigr|
    \le M\,\bigl|\kappa(\theta)-\kappa^\ast(\theta)\bigr| .
  \]
\end{lemma}

\begin{proof}

\leanok
  The two speeds share the numerator $n=1+(\min(\|z\|,R))^2\in[1,1+R^2]$ and have
  denominators $D=2\max(\kappa(\theta)-c,\delta)$, $D^\ast=2\max(\kappa^\ast(\theta)-c,
  \delta)$ with $c=\langle z,ie^{i\theta}\rangle$; both are $\ge2\delta$, and since
  $x\mapsto\max(x,\delta)$ is $1$-Lipschitz, $|D-D^\ast|\le2|\kappa(\theta)-
  \kappa^\ast(\theta)|$. Hence
  \[
    \Bigl|\frac nD-\frac n{D^\ast}\Bigr|
    =\frac{n\,|D-D^\ast|}{D\,D^\ast}
    \le\frac{(1+R^2)\cdot2|\kappa(\theta)-\kappa^\ast(\theta)|}{4\delta^2}
    =M\,|\kappa(\theta)-\kappa^\ast(\theta)| .
  \]
  (The quotient-difference bound already extracted for
  \cref{lem:truncated_speed_lipschitz} applies verbatim with the numerator difference
  set to zero.)
\end{proof}

\begin{lemma}\leanok
[combined field sensitivity]
  \label{lem:truncated_field_sub_le}
  \lean{Gluck.truncatedField_sub_le}
  \uses{def:truncated_field, lem:truncated_field_lipschitz, lem:truncated_speed_sub_le}
  Let $0\le R$, $0<\delta$, $M=\tfrac{1+R^2}{2\delta^2}$, and let $L$ be \emph{any}
  Lipschitz constant of the truncated field in $z$ (a witness of
  \cref{lem:truncated_field_lipschitz}, consumed as a hypothesis so that one fixed $L$
  can be threaded through the Gr\"onwall argument). Then for all $\kappa,\kappa^\ast$,
  $\theta$, $z$, $\tilde z$,
  \[
    \|F_{\kappa,R,\delta}(\theta,z)-F_{\kappa^\ast,R,\delta}(\theta,\tilde z)\|
    \le L\,\|z-\tilde z\|+M\,|\kappa(\theta)-\kappa^\ast(\theta)| .
  \]
\end{lemma}

\begin{proof}

\leanok
  Insert the mixed term $F_{\kappa,R,\delta}(\theta,\tilde z)$ and use the triangle
  inequality: the first difference is controlled by
  \cref{lem:truncated_field_lipschitz}; the second is
  $(\hat q_{\kappa}(\theta,\tilde z)-\hat q_{\kappa^\ast}(\theta,\tilde z))\,e^{i\theta}$,
  of norm $\le M|\kappa(\theta)-\kappa^\ast(\theta)|$ by
  \cref{lem:truncated_speed_sub_le}.
\end{proof}

\begin{lemma}\leanok
[Gr\"onwall with $L^1$ drive]
  \label{lem:gronwall_L1_drive}
  \lean{Gluck.gronwall_L1_drive}
  Let $T\ge0$, $L\ge0$, $d_0\ge0$, and let $d,g:[0,T]\to\mathbb R$ be continuous with
  $d\ge0$ and $g\ge0$ pointwise. If
  \[
    d(t)\;\le\;d_0+\int_0^t\bigl(L\,d(s)+g(s)\bigr)\,ds
    \qquad\text{for all }t\in[0,T],
  \]
  then, pointwise,
  \[
    d(t)\;\le\;e^{LT}\Bigl(d_0+\int_0^T g\Bigr)\qquad\text{for all }t\in[0,T] .
  \]
  (Project-side infrastructure: Mathlib's Gr\"onwall lemmas bound the derivative by
  $K\|f\|+\varepsilon$ with a \emph{constant} $\varepsilon$; here the drive $g$ is only
  small in $L^1$, which is exactly the regime of the Dahlberg-style reparametrization.
  The Lean statement quantifies the conclusion pointwise over $t\in[0,T]$ rather than
  through a $\sup$; the hypothesis $d\ge0$ is kept for blueprint fidelity though the
  proof does not use it.)
\end{lemma}

\begin{proof}

\leanok
  Set $u(t)=d_0+\int_0^t(L\,d+g)$. The integrand is continuous, so by the fundamental
  theorem of calculus $u$ is differentiable with $u'(t)=L\,d(t)+g(t)$, and $d\le u$ by
  hypothesis, so $u'(t)\le L\,u(t)+g(t)$ ($L\ge0$, and $u\ge d\ge0$). Consider
  $v(t)=e^{-Lt}u(t)-\int_0^t g$: it is differentiable with
  \[
    v'(t)=e^{-Lt}\bigl(u'(t)-L\,u(t)\bigr)-g(t)\le e^{-Lt}g(t)-g(t)\le0,
  \]
  using $0<e^{-Lt}\le1$ and $g\ge0$. A differentiable function with nonpositive
  derivative on an interval is antitone, so $v(t)\le v(0)=d_0$, i.e.
  $e^{-Lt}u(t)\le d_0+\int_0^t g\le d_0+\int_0^T g$ ($g\ge0$ again). Hence
  $d(t)\le u(t)\le e^{Lt}\bigl(d_0+\int_0^T g\bigr)\le e^{LT}\bigl(d_0+\int_0^T g\bigr)$.
\end{proof}

\begin{lemma}

\leanok
[invariant admissible domain — perturbative margin transport]
  \label{lem:invariant_admissible_domain}
  \lean{Gluck.invariant_admissible_domain}
  \uses{def:truncated_speed, def:spherical_flow, lem:spherical_flow_spec,
        lem:truncated_field_sub_le, lem:gronwall_L1_drive,
        lem:sphere_curvature_lower_bound}
  Let $\kappa,\kappa^\ast$ be continuous with $\kappa\ge\kappa_0>0$ pointwise, let
  $0\le R$, $0<\delta$, and let $z^\ast:[0,2\pi]\to\mathbb C$ be a trajectory of the
  $\kappa^\ast$-truncated flow that is \emph{admissible with margin} $\mu>0$:
  \[
    \sup_\theta\|z^\ast(\theta)\|\le R-\mu,
    \qquad
    \sup_\theta\,\langle z^\ast(\theta),\,i e^{i\theta}\rangle\le\kappa_0-\delta-\mu .
  \]
  Let $L$ be any Lipschitz constant of the truncated field in $z$ (a witness of
  \cref{lem:truncated_field_lipschitz}; in Lean $L:\mathbb R_{\ge0}$, coerced to
  $\mathbb R$ inside the exponential) and
  $M=\tfrac{1+R^2}{2\delta^2}$. Then every trajectory $z$ of the $\kappa$-truncated flow
  with
  \[
    e^{2\pi L}\Bigl(\|z(0)-z^\ast(0)\|+M\!\int_0^{2\pi}\!|\kappa-\kappa^\ast|\Bigr)\le\mu
  \]
  is admissible: $\|z(\theta)\|\le R$ and
  $\kappa(\theta)-\langle z(\theta),i e^{i\theta}\rangle\ge\delta$ for all
  $\theta\in[0,2\pi]$.
\end{lemma}

\begin{proof}

\leanok
  \emph{Field sensitivity.} By \cref{lem:truncated_field_sub_le},
  \[
    \|F_{\kappa}(\theta,z)-F_{\kappa^\ast}(\theta,\tilde z)\|
    \le L\|z-\tilde z\|+M\,|\kappa(\theta)-\kappa^\ast(\theta)| .
  \]
  \emph{Integral inequality.} Both trajectories are $C^1$ on $[0,2\pi]$ with derivatives
  given by their fields (\cref{lem:spherical_flow_spec}), so by the fundamental theorem
  of calculus $z(\theta)-z^\ast(\theta)=z(0)-z^\ast(0)+\int_0^\theta\bigl(F_\kappa(s,
  z(s))-F_{\kappa^\ast}(s,z^\ast(s))\bigr)ds$; taking norms and using the field
  sensitivity, $d(\theta)=\|z(\theta)-z^\ast(\theta)\|$ satisfies
  $d(\theta)\le d(0)+\int_0^\theta\bigl(L\,d+M|\kappa-\kappa^\ast|\bigr)$, with $d$
  continuous and the drive $g=M|\kappa-\kappa^\ast|$ continuous and nonnegative.
  \emph{Gr\"onwall with $L^1$ drive.} \cref{lem:gronwall_L1_drive} gives
  $\sup_{[0,2\pi]}d\le e^{2\pi L}\bigl(d(0)+M\|\kappa-\kappa^\ast\|_{L^1[0,2\pi]}\bigr)
   \le\mu$.
  \emph{Margin propagation.} Then
  $\|z(\theta)\|\le\|z^\ast(\theta)\|+\mu\le R$, and since
  $|\langle z-z^\ast,\,i e^{i\theta}\rangle|\le d(\theta)\le\mu$,
  \[
    \kappa(\theta)-\langle z,i e^{i\theta}\rangle
    \ge\kappa_0-\langle z^\ast,i e^{i\theta}\rangle-\mu\ge\delta .
  \]
  \emph{Why perturbative is the right shape.} No direct confinement estimate uniform in
  $\kappa_0$ exists: the radial velocity $d\|z\|^2/d\theta=2q\langle z,e^{i\theta}\rangle$
  changes sign, no disk is flow-invariant (at a radial critical point with
  $c=\langle z,ie^{i\theta}\rangle=-\|z\|$ one computes that a local maximum of $\|z\|$ is
  allowed exactly when $\kappa(\theta)\ge(1-\|z\|^2)/(2\|z\|)$), and for
  $\kappa_0<1/\sqrt3$ the disk of the model radius already contains denominator-vanishing
  states — admissibility is irreducibly a condition on the pair
  $(\|z\|,\langle z,ie^{i\theta}\rangle)$. The margins are instead supplied by an
  \emph{explicit} model (the symmetric step bicircle of \cref{sec:sphere_winding}, whose
  arcs give $\langle z^\ast,ie^{i\theta}\rangle=\langle w_j,ie^{i\theta}\rangle-r_j
  \le\|w_j\|-r_j$ on the $j$-th arc), and the $L^1$ distance is made arbitrarily small by
  the Dahlberg-style reparametrization of \cref{chap:dahlberg_step1}.
\end{proof}

\begin{definition}[periodic extension of a fundamental-window function]
  \label{def:periodic_extension}
  \lean{Gluck.periodicExtension}
  For $z:\mathbb R\to\mathbb C$, the \emph{$2\pi$-periodic extension} is
  $\tilde z(t)=z\bigl(t-2\pi\lfloor t/(2\pi)\rfloor\bigr)$: every $t$ is reduced by an
  integer multiple of $2\pi$ into the fundamental window $[0,2\pi)$ and $z$ is evaluated
  there. This is the device promised in \cref{lem:reconstruction_ode} for promoting a
  trajectory defined on $[0,2\pi]$ to a globally defined closed curve.
\end{definition}

\begin{lemma}[the periodic extension is periodic]
  \label{lem:periodic_extension_periodic}
  \lean{Gluck.periodicExtension_periodic}
  \uses{def:periodic_extension}
  For every $z$, the extension $\tilde z$ of \cref{def:periodic_extension} is
  $2\pi$-periodic: $\tilde z(t+2\pi)=\tilde z(t)$ for all $t$.
\end{lemma}

\begin{proof}
  $\lfloor(t+2\pi)/(2\pi)\rfloor=\lfloor t/(2\pi)\rfloor+1$, so the reduced arguments
  agree: $t+2\pi-2\pi(\lfloor t/(2\pi)\rfloor+1)=t-2\pi\lfloor t/(2\pi)\rfloor$.
\end{proof}

\begin{lemma}

\leanok
[truncation removal: admissible closed trajectories realize $\kappa$]
  \label{lem:reconstruction_ode}
  \lean{Gluck.reconstruction_ode}
  \uses{lem:truncated_speed_eq, lem:spherical_flow_spec, lem:truncated_speed_pos,
        lem:spherical_speed_solve, def:realizes_spherical_curvature,
        lem:spherical_speed_periodic, def:periodic_extension,
        lem:periodic_extension_periodic}
  Let $\kappa$ be a curvature function, $0<\delta$, $R<1$, and let
  $z:[0,2\pi]\to\mathbb C$ be a trajectory of the truncated flow that is admissible
  ($\|z(\theta)\|\le R$ and $\kappa(\theta)-\langle z(\theta),ie^{i\theta}\rangle\ge\delta$
  for all $\theta$) and closed ($z(2\pi)=z(0)$). Then the $2\pi$-periodic extension
  $\tilde z$ of $z$ (\cref{def:periodic_extension}) agrees with $z$ on $[0,2\pi]$, is a
  closed curve, and realizes $\kappa$ (\cref{def:realizes_spherical_curvature}) in the
  gauge $\varphi(\theta)=\theta$ — in particular it is $C^1$, confined to the open unit
  disk, and solves the \emph{true} reconstruction ODE $z'=q_\kappa(\theta,z)e^{i\theta}$.
  (The Lean statement packages the conclusion as the conjunction
  $\mathrm{EqOn}\wedge\mathrm{IsClosedCurve}\wedge\mathrm{RealizesSphericalCurvature}$
  for $\tilde z$; the hypothesis $0\le R$ turned out to be unnecessary and is dropped.)
\end{lemma}

\begin{proof}

\leanok
  On the admissible set $\hat q=q_\kappa$ (\cref{lem:truncated_speed_eq}), so the flow
  equation of \cref{lem:spherical_flow_spec} becomes the true ODE on $[0,2\pi]$; the
  right-hand side is continuous, so $z$ is $C^1$ there. Closedness plus
  $2\pi$-periodicity of the field ($\kappa$ periodic,
  \cref{lem:spherical_speed_periodic}) make the periodic extension a global $C^1$
  solution (one-sided derivatives match at the seam because the field values match).
  The speed $q_\kappa>0$ along the trajectory (\cref{lem:spherical_speed_solve} shape;
  positivity from admissibility), so $z'\ne0$, the tangent angle is $\varphi(\theta)=
  \theta$ itself, $\|z'\|=q_\kappa$, and the speed relation of
  \cref{def:realizes_spherical_curvature} is exactly the algebraic solve of
  \cref{def:spherical_speed}. Confinement $\|z\|\le R<1$ gives the open-disk condition.
\end{proof}

\section{Endpoint winding over the symmetric step model (S2-D)}
\label{sec:sphere_winding}

\begin{lemma}[arc bracket identity]
  \label{lem:constant_arc_inner}
  \lean{Gluck.constantArc_inner}
  For all $w\in\mathbb C$, $r,\theta\in\mathbb R$:
  $\bigl\langle w-i\,r\,e^{i\theta},\,i e^{i\theta}\bigr\rangle
   =\langle w,ie^{i\theta}\rangle-r$.
\end{lemma}

\begin{proof}
  Bilinearity of the real inner product and
  $\langle ie^{i\theta},ie^{i\theta}\rangle=\|ie^{i\theta}\|^2=1$
  (\texttt{real\_inner\_smul\_left}, \texttt{real\_inner\_self\_eq\_norm\_sq}).
\end{proof}

\begin{lemma}[arc norm expansion]
  \label{lem:constant_arc_norm_sq}
  \lean{Gluck.constantArc_norm_sq}
  For all $w\in\mathbb C$, $r,\theta\in\mathbb R$:
  $\|w-i\,r\,e^{i\theta}\|^2=\|w\|^2-2r\langle w,ie^{i\theta}\rangle+r^2$.
\end{lemma}

\begin{proof}
  The real polarization expansion $\|u-v\|^2=\|u\|^2-2\langle u,v\rangle+\|v\|^2$
  (\texttt{norm\_sub\_sq\_real}) with $\|i\,r\,e^{i\theta}\|^2=r^2$ and
  \cref{lem:constant_arc_inner}-style bilinearity for the cross term.
\end{proof}

\begin{lemma}[consistency identity of the arc data]
  \label{lem:constant_arc_consistency}
  \lean{Gluck.constantArc_consistency}
  \uses{def:spherical_speed, lem:constant_arc_norm_sq, lem:constant_arc_inner}
  Fix $K$ and start data $(\theta_0,z_0)$ with positive bracket
  $K-\langle z_0,ie^{i\theta_0}\rangle>0$, and set $r=q_K(\theta_0,z_0)$,
  $w=z_0+i\,r\,e^{i\theta_0}$. Then $1+\|w\|^2=2rK+r^2$.
\end{lemma}

\begin{proof}
  Expand $\|w\|^2=\|z_0\|^2+2r\langle z_0,ie^{i\theta_0}\rangle+r^2$
  (\texttt{norm\_add\_sq\_real}), then substitute the defining relation
  $2r\bigl(K-\langle z_0,ie^{i\theta_0}\rangle\bigr)=1+\|z_0\|^2$ of
  $r=q_K(\theta_0,z_0)$ (\cref{def:spherical_speed}).
\end{proof}

\begin{lemma}

\leanok
[constant-curvature arcs are explicit circular arcs]
  \label{lem:constant_curvature_arc}
  \lean{Gluck.constantCurvature_arc}
  \uses{def:spherical_speed, lem:spherical_flow_unique, lem:constant_arc_inner,
        lem:constant_arc_norm_sq, lem:constant_arc_consistency}
  Fix $K>0$, a start $(\theta_0,z_0)$ with $K-\langle z_0,ie^{i\theta_0}\rangle>0$, and set
  \[
    r=\frac{1+\|z_0\|^2}{2\bigl(K-\langle z_0,ie^{i\theta_0}\rangle\bigr)}=q_K(\theta_0,z_0),
    \qquad w=z_0+i\,r\,e^{i\theta_0}.
  \]
  Then $z(\theta)=w-i\,r\,e^{i\theta}$ solves the true reconstruction ODE
  $z'=q_K(\theta,z)e^{i\theta}$ with constant curvature $K$, pointwise at every $\theta$
  where the bracket stays positive ($K-\langle z(\theta),ie^{i\theta}\rangle>0$): at each
  such $\theta$ one has $q_K(\theta,z(\theta))=r$ and hence
  $z'(\theta)=q_K(\theta,z(\theta))\,e^{i\theta}$. Along the arc
  $\langle z(\theta),ie^{i\theta}\rangle=\langle w,ie^{i\theta}\rangle-r$. Moreover the
  Euclidean data satisfy the \emph{consistency identity} $K=(1+\|w\|^2-r^2)/(2r)$
  (generalizing the centered case $w=0$ of \cref{lem:spherical_speed_circle}).
\end{lemma}

\begin{proof}

\leanok
  Differentiate: $z'(\theta)=r e^{i\theta}$, so the claim is $q_K(\theta,z(\theta))=r$.
  Compute $\langle z,ie^{i\theta}\rangle=\langle w,ie^{i\theta}\rangle-r$ and
  $\|z\|^2=\|w\|^2-2r\langle w,ie^{i\theta}\rangle+r^2$; then
  $1+\|z\|^2=2r(K-\langle z,ie^{i\theta}\rangle)$ reduces to
  $1+\|w\|^2=2rK+r^2$, i.e.\ the consistency identity, which holds at $\theta_0$ by the
  definitions of $r$ and $w$ (expand $\|w\|^2=\|z_0\|^2+2r\langle z_0,ie^{i\theta_0}
  \rangle+r^2$). Spherically: stereographic images of spherical circles are Euclidean
  circles, and a constant-$K$ arc through $(\theta_0,z_0)$ is the circle of Euclidean
  radius $r=q$ tangent to $e^{i\theta_0}$ at $z_0$. Identification with the flow (where
  admissible, with clamps inactive) is \cref{lem:spherical_flow_unique}.
\end{proof}

\begin{lemma}[half-turn invariance of the truncated speed]
  \label{lem:truncated_speed_half_turn}
  \lean{Gluck.truncatedSpeed_half_turn}
  \uses{def:truncated_speed}
  If $\kappa(\theta+\pi)=\kappa(\theta)$, then
  $\hat q_{\kappa,R,\delta}(\theta+\pi,-z)=\hat q_{\kappa,R,\delta}(\theta,z)$ for all
  $\theta,z$.
\end{lemma}

\begin{proof}
  $\|{-z}\|=\|z\|$ and $ie^{i(\theta+\pi)}=-ie^{i\theta}$, so
  $\langle -z,ie^{i(\theta+\pi)}\rangle=\langle z,ie^{i\theta}\rangle$; with
  $\kappa(\theta+\pi)=\kappa(\theta)$ every ingredient of the clamped quotient is
  unchanged.
\end{proof}

\begin{lemma}[half-turn anti-equivariance of the truncated field]
  \label{lem:truncated_field_half_turn}
  \lean{Gluck.truncatedField_half_turn}
  \uses{def:truncated_field, lem:truncated_speed_half_turn}
  If $\kappa(\theta+\pi)=\kappa(\theta)$, then
  $F_{\kappa,R,\delta}(\theta+\pi,-z)=-F_{\kappa,R,\delta}(\theta,z)$ for all
  $\theta,z$.
\end{lemma}

\begin{proof}
  $F(\theta+\pi,-z)=\hat q(\theta+\pi,-z)\,e^{i(\theta+\pi)}
   =\hat q(\theta,z)\cdot(-e^{i\theta})=-F(\theta,z)$ by
  \cref{lem:truncated_speed_half_turn} and $e^{i\pi}=-1$.
\end{proof}

\begin{lemma}

\leanok
[half-turn equivariance]
  \label{lem:flow_half_turn_equivariance}
  \lean{Gluck.flow_half_turn_equivariance}
  \uses{def:truncated_speed, lem:truncated_field_solution_unique, lem:spherical_flow_spec,
        lem:truncated_speed_half_turn, lem:truncated_field_half_turn}
  If $\kappa(\theta+\pi)=\kappa(\theta)$ for all $\theta$, then the truncated field is
  invariant under the half turn: $\hat q(\theta+\pi,-z)=\hat q(\theta,z)$, hence if
  $z(\cdot)$ solves the truncated ODE on $[0,2\pi]$ then $y(\theta)=-z(\theta+\pi)$
  solves it on $[0,\pi]$. Consequently, writing $H$ for the half-period map
  $z_0\mapsto\Phi(z_0,\pi)$ of a $\pi$-periodic curvature: if $H(z_0)=-z_0$ then the
  trajectory through $z_0$ satisfies the central symmetry $z(\theta+\pi)=-z(\theta)$ for
  $\theta\in[0,\pi]$; in particular it closes at $2\pi$: $z(2\pi)=-z(\pi)=z_0$.
\end{lemma}

\begin{proof}

\leanok
  $\|{-z}\|=\|z\|$ and $\langle -z,\,ie^{i(\theta+\pi)}\rangle=\langle -z,\,-ie^{i\theta}
  \rangle=\langle z,ie^{i\theta}\rangle$, and $\kappa(\theta+\pi)=\kappa(\theta)$; so
  every ingredient of $\hat q$ is unchanged. For $\theta\in[0,\pi]$ one has
  $\theta+\pi\in[\pi,2\pi]\subseteq[0,2\pi]$, so
  $y'(\theta)=-z'(\theta+\pi)=-\hat q\,e^{i(\theta+\pi)}=\hat q(\theta,y(\theta))e^{i\theta}$
  holds on $[0,\pi]$ — $y$ solves the truncated ODE there (only there: this is why the
  comparison below happens on the half interval). If $H(z_0)=-z_0$, then
  $y(0)=-z(\pi)=z_0=z(0)$, and $z$ restricted to $[0,\pi]$ solves the same ODE
  (\cref{lem:spherical_flow_spec}), so $y=z$ on $[0,\pi]$ by the two-solution
  uniqueness on $[0,\pi]$ (\cref{lem:truncated_field_solution_unique} with $T=\pi$),
  i.e.\ $z(\theta+\pi)=-z(\theta)$ for $\theta\in[0,\pi]$; at $\theta=\pi$ this gives
  $z(2\pi)=-z(\pi)=z_0$.
\end{proof}

\subsection*{Design overview (S2-D, resolved iter-048)}

The two formerly-open sub-questions of this phase are now resolved at design level, and
the whole phase is elementary algebra plus the in-tree winding toolkit. The shape:

\textbf{DO-NOT (degeneracy of the constant model).} For \emph{constant} $\kappa=K$ every
admissible initial point lies on a spherical circle of curvature $K$
(\cref{lem:constant_curvature_arc}), so every trajectory closes and the endpoint error
map vanishes identically: the constant family carries no degree and must NOT be the
homotopy target. The fix is the \emph{symmetric step bicircle} with \emph{small
contrast}: levels $a=c-h/2$, $b=c+h/2$ around the midpoint $c$ of the four-vertex
overlap window, on the in-tree canonical equal-quarter breakpoints $0,\pi/2,\pi,3\pi/2$
(\cref{def:four_arc_step_function} — note $\kappa^\ast=$
$\texttt{stepCurvature } b\,a\,0\,(\pi/2)\,\pi\,(3\pi/2)$ takes the value $a$ on
$[0,\pi/2)\cup[\pi,3\pi/2)$ and $b$ on the complement, and is $\pi$-periodic).

\textbf{Everything is explicit arc algebra.} Because the step is piecewise constant, its
model trajectories are concatenations of \emph{explicit circular arcs}
(\cref{lem:constant_curvature_arc}); no flow, no variational equation, and no
$\mathrm{PicardLindel\ddot of}$ package is needed on the model side. The model endpoint
error map $E^\ast_{a,b}$ (\cref{def:step_error_map}) is a closed-form composition of
four arc maps. Its degeneracy at $h=0$ becomes a feature: $E^\ast_{c,c}\equiv0$
identically, so $E^\ast_{a,b}=h\,W+O(h^2)$ for an explicit \emph{first-variation field}
$W$, and the planner's symbolic computation (numerically validated: the arc composite
matches an RK4 integration of the step ODE to $10^{-7}$, $E^\ast(z_0^\ast)=O(h^2)$, and
$E^\ast/h$ matches the predicted linearization with ratio $\to1$) gives
\[
  W(z_0^\ast)=0,\qquad
  DW(z_0^\ast)=-\tfrac{2r^\ast}{1+c^2}\cdot\overline{(\,\cdot\,)}
  \qquad(z_0^\ast=-i\,r^\ast,\ r^\ast=\sqrt{1+c^2}-c),
\]
a \emph{conjugation-type} (orientation-reversing) nondegenerate linear map: the model
error loop on a small circle $\partial B(z_0^\ast,\rho)$ has winding $-1\ne0$. The exact
mechanism behind the cancellations is the \emph{quadratic identity}
\cref{lem:speed_quadratic_identity}: $q_c(\theta,z)-r^\ast$ vanishes to second order in
the distance to the centered circle, with an exact rational formula (no Taylor
remainders anywhere: all expansions in this section are exact algebraic identities plus
norm estimates of explicitly-controlled factors).

\textbf{Transport.} The comparison of the actual (reparametrized) $\kappa$-flow with the
step model cannot use \cref{lem:invariant_admissible_domain} directly ($\kappa^\ast$ is
discontinuous, and the concatenated model has corners), and Mathlib-style continuity of
the drive fails at the breakpoints. The resolution is \emph{arcwise chaining}: on each
quarter arc the model level is a \emph{constant} $K\in\{a,b\}$, the drive
$M|\kappa-K|$ is continuous there, and the single-arc transport
(\cref{lem:invariant_admissible_arc}) is the verbatim margin-transport proof on a
shifted interval. Four chained applications give the full-period transport
(\cref{lem:step_model_transport}) with constant $e^{2\pi L}$, both for admissibility and
for the endpoint estimate $\|E_\kappa(z_0)-E^\ast(z_0)\|\le e^{2\pi L}M\!\int|\kappa-
\kappa^\ast|$.

\textbf{Quantifier order.} $c$ and $\kappa_0=\min\kappa$ are fixed by $\kappa$; margins
$(R,\delta,\mu,\rho_0,h_0)$ by \cref{lem:step_model_margins}; the disk radius $\rho$ and
contrast $h$ by \cref{lem:step_error_expansion} (boundary margin
$\sim\tfrac{r^\ast}{1+c^2}h\rho$); the reparametrization quality $\varepsilon\ll h\rho$
last (it is free, \cref{lem:step_L1_reparam}). The Euclidean
\texttt{errorMap\_winding\_eq\_one} \emph{value} is not reusable; only
\cref{thm:existence_of_zero} and the perturbation/congruence lemmas of
\cref{chap:winding} are.

\begin{definition}[spherical arc map]
  \label{def:spherical_arc_map}
  \lean{Gluck.sphericalArcMap}
  \uses{def:spherical_speed}
  For $K,\theta_0,\Delta\in\mathbb R$ and $z\in\mathbb C$, the \emph{arc map} is
  \[
    A_{K,\theta_0,\Delta}(z)\;=\;z+i\,q_K(\theta_0,z)\,e^{i\theta_0}\bigl(1-e^{i\Delta}\bigr),
  \]
  where $q_K$ abbreviates the gauge speed $q_{(\theta\mapsto K)}$ of
  \cref{def:spherical_speed} for the constant curvature $K$. By
  \cref{lem:constant_curvature_arc}, where the bracket stays positive this is the
  time-$\Delta$ endpoint of the constant-$K$ trajectory started at $(\theta_0,z)$: the
  arc trajectory is $\theta\mapsto w-i\,r\,e^{i\theta}$ with $r=q_K(\theta_0,z)$,
  $w=z+i\,r\,e^{i\theta_0}$, and its value at $\theta_0+\Delta$ is exactly
  $A_{K,\theta_0,\Delta}(z)$. (A total function of the junk-value kind, like $q$ itself:
  the geometric meaning requires bracket positivity, the definition does not.)
\end{definition}

\begin{lemma}[quadratic identity: exact second-order vanishing of $q$ at the centered circle]
  \label{lem:speed_quadratic_identity}
  \lean{Gluck.sphericalSpeed_sub_radius}
  \uses{def:spherical_speed, lem:spherical_speed_circle}
  Fix $c>0$ and set $r^\ast=\sqrt{1+c^2}-c$, so that $r^{\ast2}+2cr^\ast-1=0$. For all
  $\theta$ and all $z$ with $D:=c-\langle z,ie^{i\theta}\rangle\ne0$,
  \[
    q_c(\theta,z)-r^\ast\;=\;\frac{\bigl\|z+r^\ast i e^{i\theta}\bigr\|^2}{2D}
    \;=\;\frac{\|z-z_c(\theta)\|^2}{2D},
    \qquad z_c(\theta):=-r^\ast i e^{i\theta}
  \]
  ($z_c$ is the centered circle of \cref{lem:spherical_speed_circle}). In particular
  $q_c-r^\ast$ vanishes \emph{exactly to second order} in the distance to the centered
  circle: the $z$-gradient of $q_c(\theta,\cdot)$ vanishes at $z_c(\theta)$, and on
  $D>0$ one has $q_c\ge r^\ast$.
\end{lemma}

\begin{proof}
  \uses{lem:constant_arc_norm_sq, lem:constant_arc_inner}
  $q_c-r^\ast=\dfrac{(1+\|z\|^2)-2r^\ast D}{2D}$, and the numerator is
  \[
    1+\|z\|^2-2r^\ast c+2r^\ast\langle z,ie^{i\theta}\rangle
    =\|z\|^2+r^{\ast2}+2\langle z,\,r^\ast ie^{i\theta}\rangle
    =\|z+r^\ast ie^{i\theta}\|^2,
  \]
  using $1-2r^\ast c=r^{\ast2}$ (the defining identity of $r^\ast$) and the polarization
  expansion (\cref{lem:constant_arc_norm_sq} with $-r$ replaced by $r^\ast$).
\end{proof}

\begin{lemma}[speed half-turn (true speed)]
  \label{lem:spherical_speed_half_turn}
  \lean{Gluck.sphericalSpeed_half_turn}
  \uses{def:spherical_speed}
  For every constant $K$: $q_K(\theta+\pi,-z)=q_K(\theta,z)$ for all $\theta,z$. (Same
  computation as \cref{lem:truncated_speed_half_turn}, without clamps: $\|-z\|=\|z\|$ and
  $\langle -z,ie^{i(\theta+\pi)}\rangle=\langle z,ie^{i\theta}\rangle$.)
\end{lemma}

\begin{proof}
  \uses{lem:truncated_speed_half_turn}
  As stated; $ie^{i(\theta+\pi)}=-ie^{i\theta}$ and bilinearity of the inner product.
\end{proof}

\begin{lemma}[arc-map half-turn anti-equivariance]
  \label{lem:arc_map_half_turn}
  \lean{Gluck.sphericalArcMap_half_turn}
  \uses{def:spherical_arc_map, lem:spherical_speed_half_turn}
  $A_{K,\theta_0+\pi,\Delta}(-z)=-A_{K,\theta_0,\Delta}(z)$ for all
  $K,\theta_0,\Delta,z$.
\end{lemma}

\begin{proof}
  $q_K(\theta_0+\pi,-z)=q_K(\theta_0,z)$ (\cref{lem:spherical_speed_half_turn}) and
  $e^{i(\theta_0+\pi)}=-e^{i\theta_0}$, so every term of
  \cref{def:spherical_arc_map} flips sign.
\end{proof}

\begin{lemma}[arc concatenation]
  \label{lem:arc_map_concat}
  \lean{Gluck.sphericalArcMap_concat}
  \uses{def:spherical_arc_map, lem:constant_curvature_arc}
  Fix $K>0$, start data $(\theta_0,z)$ with positive bracket, and $\Delta_1,\Delta_2\ge0$.
  If the bracket $K-\langle z(\theta),ie^{i\theta}\rangle$ stays positive along the arc
  trajectory $z(\theta)=w-ire^{i\theta}$ for $\theta\in[\theta_0,\theta_0+\Delta_1+
  \Delta_2]$, then
  \[
    A_{K,\theta_0+\Delta_1,\Delta_2}\bigl(A_{K,\theta_0,\Delta_1}(z)\bigr)
    =A_{K,\theta_0,\Delta_1+\Delta_2}(z) .
  \]
\end{lemma}

\begin{proof}
  By \cref{lem:constant_curvature_arc}(iii), $q_K$ is constant $=r$ along the arc where
  the bracket is positive; in particular
  $q_K(\theta_0+\Delta_1,\,z(\theta_0+\Delta_1))=r=q_K(\theta_0,z)$, so the second arc
  continues the same circle: both sides equal $w-ire^{i(\theta_0+\Delta_1+\Delta_2)}$
  after expanding \cref{def:spherical_arc_map} and collecting the exponentials.
  In particular the full-turn map is the identity where admissible:
  $A_{K,\theta_0,2\pi}(z)=z$ (since $e^{2\pi i}=1$), which is the exact form of the
  constant-model degeneracy.
\end{proof}

\begin{definition}[step half-period map]
  \label{def:step_half_map}
  \lean{Gluck.stepHalfMap}
  \uses{def:spherical_arc_map}
  For levels $a,b$, the \emph{half-period map} of the symmetric equal-quarter step is
  the composition of the two quarter-arc maps
  \[
    H_{a,b}(z)\;=\;A_{b,\pi/2,\pi/2}\bigl(A_{a,0,\pi/2}(z)\bigr) .
  \]
  (First quarter has level $a$, second $b$, matching
  $\kappa^\ast=\texttt{stepCurvature } b\,a\,0\,(\pi/2)\,\pi\,(3\pi/2)$ of
  \cref{def:four_arc_step_function}.)
\end{definition}

\begin{definition}[step model error map]
  \label{def:step_error_map}
  \lean{Gluck.stepErrorMap}
  \uses{def:step_half_map}
  $E^\ast_{a,b}(z)\;=\;-H_{a,b}\bigl(-H_{a,b}(z)\bigr)-z$. By
  \cref{lem:arc_map_half_turn}, the arcs of the second half period (levels $a,b$ at
  $\theta_0=\pi,3\pi/2$, by $\pi$-periodicity of the step) are obtained from the first
  half by the half turn, so $z+E^\ast_{a,b}(z)$ is exactly the four-arc composite
  \[
    A_{b,3\pi/2,\pi/2}\circ A_{a,\pi,\pi/2}\circ A_{b,\pi/2,\pi/2}\circ A_{a,0,\pi/2}
  \]
  — the endpoint at $2\pi$ of the concatenated step-model trajectory. At $a=b=c$,
  $E^\ast_{c,c}(z)=0$ wherever the circle through $z$ stays bracket-positive
  (\cref{lem:arc_map_concat}).
\end{definition}

\begin{lemma}[explicit arcs solve the truncated ODE]
  \label{lem:constant_arc_solves_truncated}
  \lean{Gluck.constantArc_solves_truncated}
  \uses{lem:constant_curvature_arc, lem:truncated_speed_eq, def:truncated_field}
  Fix a constant $K$, $0<\delta$, $R$, start data $(\theta_0,z_0)$, and let
  $z(\theta)=w-ire^{i\theta}$ be the arc trajectory of \cref{lem:constant_curvature_arc}.
  If along $\theta\in[t_1,t_2]$ the arc is \emph{admissible with clamps inactive} —
  $\|z(\theta)\|\le R$ and $K-\langle z(\theta),ie^{i\theta}\rangle\ge\delta$ — then $z$
  solves the \emph{truncated} ODE for the constant curvature $K$ on $[t_1,t_2]$ in the
  $\mathrm{HasDerivWithinAt}$ sense:
  $z'(\theta)=\hat q_{K,R,\delta}(\theta,z(\theta))\,e^{i\theta}$ on $[t_1,t_2]$.
\end{lemma}

\begin{proof}
  On the admissible set the clamps are inactive, $\hat q=q_K$
  (\cref{lem:truncated_speed_eq}), and $z'(\theta)=re^{i\theta}
  =q_K(\theta,z(\theta))e^{i\theta}$ by \cref{lem:constant_curvature_arc}.
\end{proof}

\begin{lemma}[single-arc margin transport (shifted interval, constant model level)]
  \label{lem:invariant_admissible_arc}
  \lean{Gluck.invariant_admissible_arc}
  \uses{lem:truncated_field_sub_le, lem:gronwall_L1_drive,
        lem:invariant_admissible_domain}
  Let $\kappa$ be continuous with $\kappa\ge\kappa_0$ pointwise, $0<\delta$, let $L$ be a
  Lipschitz witness of \cref{lem:truncated_field_lipschitz} and
  $M=\tfrac{1+R^2}{2\delta^2}$. Let $[t_1,t_2]\subseteq\mathbb R$, $K\in\mathbb R$
  (\emph{constant} model level), and let $z,z^\ast:[t_1,t_2]\to\mathbb C$ be continuous
  with, in the $\mathrm{HasDerivWithinAt}$ sense on $[t_1,t_2]$,
  \[
    z'=\hat q_{\kappa,R,\delta}(\theta,z)\,e^{i\theta},\qquad
    z^{\ast\prime}=\hat q_{K,R,\delta}(\theta,z^\ast)\,e^{i\theta},
  \]
  and suppose the model has margins $\mu>0$ on the arc:
  $\|z^\ast(\theta)\|\le R-\mu$ and
  $\langle z^\ast(\theta),ie^{i\theta}\rangle\le\kappa_0-\delta-\mu$ for
  $\theta\in[t_1,t_2]$. If
  \[
    e^{L(t_2-t_1)}\Bigl(\|z(t_1)-z^\ast(t_1)\|
      +M\!\int_{t_1}^{t_2}\!|\kappa(s)-K|\,ds\Bigr)\;\le\;\mu,
  \]
  then for all $\theta\in[t_1,t_2]$:
  \[
    \|z(\theta)-z^\ast(\theta)\|\le
    e^{L(t_2-t_1)}\Bigl(\|z(t_1)-z^\ast(t_1)\|+M\!\int_{t_1}^{t_2}\!|\kappa-K|\Bigr),
    \quad
    \|z(\theta)\|\le R,\quad
    \kappa(\theta)-\langle z(\theta),ie^{i\theta}\rangle\ge\delta .
  \]
\end{lemma}

\begin{proof}
  Verbatim the proof of \cref{lem:invariant_admissible_domain} after the time shift
  $s\mapsto t_1+s$ (which turns $[t_1,t_2]$ into $[0,T]$, $T=t_2-t_1$, and shifts the
  fields accordingly): FTC difference identity, the field sensitivity
  \cref{lem:truncated_field_sub_le} with $\kappa^\ast\equiv K$ — the drive
  $g=M|\kappa-K|$ is \emph{continuous} because the model level is constant, which is the
  entire point of the arcwise formulation — then \cref{lem:gronwall_L1_drive} and margin
  propagation via $|\langle z-z^\ast,ie^{i\theta}\rangle|\le\|z-z^\ast\|$.
\end{proof}

\begin{lemma}[four-arc chained transport against the step model]
  \label{lem:step_model_transport}
  \lean{Gluck.stepModel_transport}
  \uses{lem:invariant_admissible_arc, lem:constant_arc_solves_truncated,
        def:step_half_map, def:step_error_map, lem:arc_map_half_turn,
        def:four_arc_step_function, def:spherical_flow, lem:spherical_flow_spec}
  Let $\kappa$ be continuous with $\kappa\ge\kappa_0$ pointwise, $0<\delta$, $L$, $M$ as
  above, levels $a,b$, and $\kappa^\ast=\texttt{stepCurvature } b\,a\,0\,(\pi/2)\,\pi\,
  (3\pi/2)$. Let $z$ be a trajectory of the $\kappa$-truncated flow on $[0,2\pi]$ with
  $z(0)=z_0$, and let the \emph{concatenated step model} from $z_0$ be the four
  quarter-arc trajectories with levels $a,b,a,b$ at $\theta_0=0,\pi/2,\pi,3\pi/2$, each
  arc started at the endpoint of the previous one. Assume each of the four arcs is
  admissible with margins $\mu$ as in \cref{lem:invariant_admissible_arc} (with its own
  constant level), and
  \[
    e^{2\pi L}\,M\!\int_0^{2\pi}\!|\kappa-\kappa^\ast|\;\le\;\mu .
  \]
  Then $z$ is admissible on $[0,2\pi]$ ($\|z\|\le R$, bracket $\ge\delta$), and the
  endpoint satisfies
  \[
    \bigl\|\,\bigl(z(2\pi)-z_0\bigr)-E^\ast_{a,b}(z_0)\,\bigr\|
    \;\le\;e^{2\pi L}\,M\!\int_0^{2\pi}\!|\kappa-\kappa^\ast| .
  \]
\end{lemma}

\begin{proof}
  Apply \cref{lem:invariant_admissible_arc} on each quarter $[\theta_j,\theta_{j+1}]$,
  $\theta_j=j\pi/2$, with the constant level $K_j\in\{a,b\}$ of that quarter; the model
  arcs solve the truncated ODE by \cref{lem:constant_arc_solves_truncated}, and on each
  quarter $\int_{\theta_j}^{\theta_{j+1}}|\kappa-K_j|=\int_{\theta_j}^{\theta_{j+1}}
  |\kappa-\kappa^\ast|$ ($\kappa^\ast=K_j$ off the single breakpoint). Chain the distance
  bounds: with $d_j=\|z(\theta_j)-z^\ast(\theta_j)\|$ and $I_j=M\int_{\theta_j}^{
  \theta_{j+1}}|\kappa-\kappa^\ast|$, each quarter gives $d_{j+1}\le e^{L\pi/2}(d_j+I_j)$
  and $d_0=0$, whence inductively $d(\theta)\le e^{2\pi L}\sum_j I_j$ throughout —
  in particular each quarter's smallness hypothesis holds and admissibility propagates
  across all four quarters. The model endpoint at $2\pi$ is the four-arc composite,
  which equals $z_0+E^\ast_{a,b}(z_0)$ by \cref{def:step_error_map} (via
  \cref{lem:arc_map_half_turn}), giving the endpoint estimate. (The prover may extract
  the composite-equals-$E^\ast$ step as a helper lemma.)
\end{proof}

\begin{lemma}[uniform margins of the step model near the centered circle]
  \label{lem:step_model_margins}
  \lean{Gluck.stepModel_margins}
  \uses{def:spherical_arc_map, lem:constant_arc_inner, lem:speed_quadratic_identity,
        lem:spherical_speed_circle, def:step_half_map}
  Fix $c>0$ and $\kappa_0\in(0,c]$, and set $r^\ast=\sqrt{1+c^2}-c$,
  $z_0^\ast=-i\,r^\ast$. There exist explicit $R<1$, $\delta>0$, $\mu>0$, $\rho_0>0$,
  $h_0>0$ (functions of $c,\kappa_0$ only) such that for all levels $a,b$ with
  $|a-c|\le h_0$, $|b-c|\le h_0$ and every $z_0$ with $\|z_0-z_0^\ast\|\le\rho_0$, each
  of the four quarter arcs of the concatenated step model from $z_0$
  (\cref{lem:step_model_transport}) satisfies, along the whole arc,
  \[
    \|z^\ast(\theta)\|\le R-\mu,\qquad
    \langle z^\ast(\theta),ie^{i\theta}\rangle\le\kappa_0-\delta-\mu,\qquad
    K_j-\langle z^\ast(\theta),ie^{i\theta}\rangle\ge\delta
  \]
  (the last condition keeps the clamps inactive so
  \cref{lem:constant_arc_solves_truncated} applies).
\end{lemma}

\begin{proof}
  Along the centered circle itself ($a=b=c$, $z_0=z_0^\ast$) one has
  $\|z^\ast\|\equiv r^\ast<1$ and $\langle z^\ast(\theta),ie^{i\theta}\rangle\equiv
  -r^\ast<0$ (\cref{lem:spherical_speed_circle}), so all three quantities have strict
  slack: e.g.\ $R=(1+r^\ast)/2$, $\delta+\mu\le\min(\kappa_0+r^\ast,c+r^\ast)/2$.
  Each arc map $A_{K,\theta_0,\pi/2}$ is explicitly Lipschitz in $(K,z)$ on the compact
  neighborhood $\{|K-c|\le c/2,\ \|z-z_c(\theta_0)\|\le r^\ast/2\}$ (the bracket is
  bounded below there via \cref{lem:constant_arc_inner} and Cauchy–Schwarz, so $q_K$ is a
  quotient of Lipschitz functions with denominator bounded away from $0$), and along each
  arc $\langle z^\ast(\theta),ie^{i\theta}\rangle=\langle w_j,ie^{i\theta}\rangle-r_j$
  with $(w_j,r_j)$ Lipschitz-close to $(0,r^\ast)$. Chaining the four Lipschitz maps,
  the whole concatenated trajectory deviates from the centered circle by at most
  $C(c)\,(\rho_0+h_0)$; choose $\rho_0,h_0$ so this is below half the slack. All
  constants are explicit; no compactness argument is needed.
\end{proof}

\begin{lemma}[first-variation expansion of the step error map]
  \label{lem:step_error_expansion}
  \lean{Gluck.stepError_expansion}
  \uses{def:step_error_map, def:step_half_map, def:spherical_arc_map,
        lem:speed_quadratic_identity, lem:arc_map_concat, lem:constant_arc_inner}
  Fix $c>0$; let $r^\ast,z_0^\ast$ be as above and set $\eta=\tfrac{2r^\ast}{1+c^2}>0$.
  There exist explicit $\rho_1>0$, $h_1>0$, $C>0$ (functions of $c$ only) such that for
  all $h\in(0,h_1]$, with levels $a=c-h/2$, $b=c+h/2$, and all $z_0$ with
  $\|z_0-z_0^\ast\|\le\rho_1$:
  \[
    \Bigl\|\,E^\ast_{a,b}(z_0)\;+\;\eta\,h\,\overline{(z_0-z_0^\ast)}\,\Bigr\|
    \;\le\;C\,h\,\bigl(\|z_0-z_0^\ast\|^2+h\bigr) .
  \]
  \emph{Provenance.} The linear part $-\eta\,h\,\overline{(\,\cdot\,)}$ is
  planner-derived (iter-048) and numerically validated ($E^\ast/h$ against the predicted
  field at multiple base points and directions, ratio $\to1$ as the offset shrinks);
  the prover should re-derive the coefficient symbolically during formalization and may
  correct its exact value — the \emph{load-bearing} content is only that the linear part
  is a nonzero real multiple of complex conjugation (orientation-reversing,
  nondegenerate).
\end{lemma}

\begin{proof}
  All expansions are exact algebra; no abstract Taylor theorem is used. Write
  $\delta z=z_0-z_0^\ast$. The two mechanisms:
  (i) \emph{level sensitivity is exact}: for any levels $K,K'$ and admissible $z$,
  \[
    q_K(\theta,z)-q_{K'}(\theta,z)
    =\frac{(1+\|z\|^2)\,(K'-K)}{2\,D_K D_{K'}},\qquad
    D_K=K-\langle z,ie^{i\theta}\rangle,
  \]
  a one-line quotient identity; at the circle base point this evaluates to
  $\mp\tfrac{h}{2}\cdot\tfrac{r^\ast}{\sqrt{1+c^2}}+O(h(h+\|\delta z\|))$ for
  $K=c\pm h/2$.
  (ii) \emph{base-point sensitivity is exactly quadratic}:
  \cref{lem:speed_quadratic_identity} gives $q_c(\theta,z)-r^\ast=\|z-z_c(\theta)\|^2/
  (2D)$, so all first-order-in-$\delta z$ contributions at level $c$ vanish and
  re-enter only through the $h$-dependent terms.
  Expand the four-arc composite $z_0+E^\ast_{a,b}(z_0)$ arc by arc using (i), (ii) and
  the gradient formula for $q_c$ implied by (ii) (each arc's output is the input plus
  $i\,q\,e^{i\theta_0}(1-e^{i\pi/2})$ with $1-e^{i\pi/2}=1-i$, and intermediate points
  stay within $O(\|\delta z\|+h)$ of the centered circle by
  \cref{lem:step_model_margins}-style bounds); after cancellation of the $h$-free part
  ($E^\ast_{c,c}\equiv0$, \cref{lem:arc_map_concat}) the coefficient of $h$ collects to
  $-\eta\,\overline{\delta z}+O(\|\delta z\|^2)$ and every remaining term carries
  $h(\|\delta z\|^2+h)$. The bookkeeping is finite (eight $q$-evaluations), and the
  prover should organize it through named intermediate \texttt{have}-identities rather
  than one monolithic \texttt{ring} call.
\end{proof}

\begin{lemma}[winding of a conjugate loop]
  \label{lem:winding_conj_loop}
  \lean{Gluck.windingNumberC_conj_loop}
  \uses{def:winding_number_c, lem:winding_number_eq_div_of_lift}
  For $w\ne0$, the loop $t\mapsto w\cdot\overline{e^{2\pi it}}$ on $[0,1]$ is nowhere
  zero and has $\mathbb C$-winding number $-1$.
\end{lemma}

\begin{proof}
  The normalization of the loop is $t\mapsto\mathrm{circleProj}(w)\cdot e^{-2\pi it}$,
  which lifts through $\mathrm{Circle.exp}$ via $t\mapsto\arg(w)-2\pi t$; by
  \cref{lem:winding_number_eq_div_of_lift} the winding number is
  $\bigl((\arg w-2\pi)-\arg w\bigr)/(2\pi)=-1$. (CAUTION — the \emph{entire} lift layer
  of \texttt{Winding.lean} is \texttt{private} and not importable: not just
  \texttt{windingNumber\_negStandard} but also \texttt{windingNumber\_eq\_div\_of\_lift},
  \texttt{angleLift}, and the raw \texttt{windingNumber}; the public surface is only
  \texttt{windingNumberC}, \texttt{windingNumberC\_congr},
  \texttt{windingNumberC\_eq\_of\_perturb}, \texttt{windingNumberC\_posScalarField},
  \texttt{diskBoundaryLoop}, and \texttt{thm:existence\_of\_zero}. The local proof must
  therefore redo the angle-lift computation self-contained from Mathlib's
  \texttt{Circle.isCoveringMap\_exp} path lifting — the same route as
  \texttt{Winding.lean}'s construction, an estimated 80--120 LOC helper layer, NOT a
  five-line replication.)
\end{proof}

\begin{lemma}[four values at freely chosen levels]
  \label{lem:exists_abab_levels}
  \lean{Gluck.exists_abab_levels}
  \uses{def:four_vertex_condition, lem:ivt_hits, lem:abab_values}
  Let $\kappa$ be continuous and $2\pi$-periodic with alternating extrema data
  $p_1<q_1<p_2<q_2<p_1+2\pi$ satisfying the strict value separation
  $\max(\kappa(q_1),\kappa(q_2))<\min(\kappa(p_1),\kappa(p_2))$. Then for \emph{every}
  pair of levels $a,b$ with
  $\max(\kappa(q_1),\kappa(q_2))<a<b<\min(\kappa(p_1),\kappa(p_2))$ there exist
  $\theta_1<\theta_2<\theta_3<\theta_4<\theta_1+2\pi$ with
  $\kappa(\theta_1,\theta_2,\theta_3,\theta_4)=(a,b,a,b)$. (This refines
  \cref{lem:abab_values}, which produces one specific pair of levels — the endpoints of
  the overlap window; here the levels are free, which is what allows the small-contrast
  choice $a=c-h/2$, $b=c+h/2$.)
\end{lemma}

\begin{proof}
  Four applications of the packaged IVT (\cref{lem:ivt_hits}): on $[q_1,p_2]$, $\kappa$
  runs from $\le a$ to $\ge b$, giving $\theta_1$ with $\kappa=a$ and then, on
  $[\theta_1,p_2]$, $\theta_2\ge\theta_1$ with $\kappa=b$ (strict inequality since the
  values differ); on $[p_2,q_2]$ ($\ge b\to\le a$) a point $\theta_3$ with $\kappa=a$;
  on $[q_2,p_1+2\pi]$ ($\le a\to\kappa(p_1)\ge b$, using periodicity) a point
  $\theta_4$ with $\kappa=b$. The ordering and window conditions follow since
  $\theta_4\le p_1+2\pi<q_1+2\pi\le\theta_1+2\pi$. (Home: this lemma lives in the S²
  file — the Euclidean files are frozen.)
\end{proof}

\begin{lemma}[$L^1$ step reparametrization]
  \label{lem:step_L1_reparam}
  \lean{Gluck.exists_step_L1_reparam}
  \uses{lem:exists_preliminary_reparam, def:four_arc_step_function,
        def:is_curvature_function}
  Let $\kappa$ be a curvature function, $0<a<b$, and $\theta_1<\theta_2<\theta_3<
  \theta_4<\theta_1+2\pi$ with $\kappa(\theta_1,\theta_2,\theta_3,\theta_4)=(a,b,a,b)$.
  For every $\varepsilon>0$ there is an orientation-preserving circle reparametrization
  $h_1$ ($\mathrm{StrictMono}$, $C^1$ with continuous positive derivative,
  $h_1(\theta+2\pi)=h_1(\theta)+2\pi$) such that
  \[
    \int_0^{2\pi}\bigl|\kappa(h_1(\theta))-\kappa^\ast(\theta)\bigr|\,d\theta
    \;<\;\varepsilon,
    \qquad \kappa^\ast=\texttt{stepCurvature } b\,a\,0\,(\pi/2)\,\pi\,(3\pi/2).
  \]
\end{lemma}

\begin{proof}
  Apply \cref{lem:exists_preliminary_reparam} with
  $\varepsilon'=\varepsilon/\bigl(B+2\pi+1\bigr)$ where
  $B=\sup_{[0,2\pi]}\kappa+b$ (finite by compactness and periodicity, and a pointwise
  bound for $|\kappa\circ h_1-\kappa^\ast|$ since $0<\kappa^\ast\le b$ and $\kappa\circ
  h_1$ ranges over the values of $\kappa$). Split the integral over the bad set (measure
  $<\varepsilon'$, integrand $\le B$) and its complement in $[0,2\pi)$ (integrand
  $\le\varepsilon'$, measure $\le2\pi$):
  $\int\le\varepsilon' B+2\pi\varepsilon'<\varepsilon$.
\end{proof}

\begin{lemma}[spherical endpoint winding: closed admissible trajectory for the
reparametrized curvature]
  \label{lem:spherical_endpoint_winding}
  \lean{Gluck.spherical_endpoint_winding}
  \uses{def:spherical_endpoint, lem:spherical_endpoint_continuousOn,
        def:step_error_map, lem:step_model_transport, lem:step_model_margins,
        lem:step_error_expansion, lem:winding_conj_loop,
        lem:winding_number_c_perturb, lem:exists_abab_levels, lem:step_L1_reparam,
        thm:existence_of_zero, def:sphere_four_vertex,
        lem:sphere_curvature_lower_bound}
  Let $\kappa$ satisfy \cref{def:sphere_four_vertex} and be non-constant. Then there
  exist $R<1$, $\delta>0$, an orientation-preserving circle reparametrization $h_1$
  ($\mathrm{StrictMono}$, $C^1$ with continuous positive derivative,
  $h_1(\theta+2\pi)=h_1(\theta)+2\pi$), and $z_0\in\mathbb C$, such that the trajectory
  of the $(\kappa\circ h_1,R,\delta)$-truncated flow through $z_0$ is \emph{closed}
  ($\Phi(z_0,2\pi)=z_0$, i.e.\ the endpoint map of \cref{def:spherical_endpoint}
  vanishes at $z_0$) and \emph{admissible} on $[0,2\pi]$:
  $\|z(\theta)\|\le R$ and
  $(\kappa\circ h_1)(\theta)-\langle z(\theta),ie^{i\theta}\rangle\ge\delta$.
\end{lemma}

\begin{proof}
  \emph{Data.} From the non-constant four-vertex branch take the alternating extrema and
  the overlap window; let $c$ be its midpoint and $\kappa_0=\min_{[0,2\pi]}\kappa>0$
  ($\le c$, shrinking $\kappa_0$ if necessary). \cref{lem:step_model_margins} for
  $(c,\kappa_0)$ fixes $R,\delta,\mu,\rho_0,h_0$; \cref{lem:step_error_expansion} fixes
  $\rho_1,h_1',C,\eta$. Choose
  $\rho=\min\bigl(\rho_0,\rho_1,\tfrac{\eta}{4C}\bigr)$ and then
  $h=\min\bigl(h_0,h_1',\tfrac{\eta\rho}{4C},\tfrac{w}{2}\bigr)$ where $w$ is the
  half-width of the overlap window; set $a=c-h/2$, $b=c+h/2$.
  \emph{Boundary margin of the model.} On the circle $\|z_0-z_0^\ast\|=\rho$,
  \cref{lem:step_error_expansion} gives
  $\|E^\ast_{a,b}(z_0)+\eta h\overline{(z_0-z_0^\ast)}\|\le Ch(\rho^2+h)\le
  \tfrac{\eta h\rho}{2}$, so $\|E^\ast_{a,b}\|\ge\tfrac{\eta h\rho}{2}$ there.
  \emph{Reparametrization.} \cref{lem:exists_abab_levels} supplies the $(a,b,a,b)$
  crossing data; \cref{lem:step_L1_reparam} then supplies $h_1$ with
  $\int_0^{2\pi}|\kappa\circ h_1-\kappa^\ast|<\varepsilon$, where $\varepsilon>0$ is
  chosen with $e^{2\pi L}M\varepsilon\le\min\bigl(\mu,\tfrac{\eta h\rho}{8}\bigr)$
  ($L$ a Lipschitz witness for the truncated field, $M=\tfrac{1+R^2}{2\delta^2}$). Note
  $\kappa\circ h_1$ is continuous, $2\pi$-periodic, and $\ge\kappa_0$.
  \emph{Winding.} For $z_0$ on the boundary circle, \cref{lem:step_model_transport}
  (margins from \cref{lem:step_model_margins}) bounds
  $\|E_{\kappa\circ h_1}(z_0)-E^\ast_{a,b}(z_0)\|\le e^{2\pi L}M\varepsilon\le
  \tfrac{\eta h\rho}{8}$, where $E_{\kappa\circ h_1}$ is the truncated endpoint map of
  \cref{def:spherical_endpoint}. Combining with the boundary margin and the expansion,
  the loop of $E_{\kappa\circ h_1}$ around $\partial B(z_0^\ast,\rho)$ is a
  perturbation, of size $<$ the model's norm, of the conjugate-linear loop
  $t\mapsto-\eta h\,\rho\,\overline{e^{2\pi it}}$; by
  \cref{lem:winding_number_c_perturb} and \cref{lem:winding_conj_loop} its winding is
  $-1\ne0$, and it is nowhere zero on the boundary.
  \emph{Zero.} $E_{\kappa\circ h_1}$ is continuous on the closed ball
  (\cref{lem:spherical_endpoint_continuousOn}); after the affine chart
  $w\mapsto z_0^\ast+\rho w$ of the unit disk, \cref{thm:existence_of_zero} yields an
  interior $z_0$ with $E_{\kappa\circ h_1}(z_0)=0$: the trajectory closes.
  \emph{Admissibility.} Apply \cref{lem:step_model_transport} once more at this interior
  $z_0$ (its model margins hold since $\|z_0-z_0^\ast\|<\rho\le\rho_0$), with the same
  $\varepsilon$-smallness: the closed trajectory is admissible on $[0,2\pi]$.
\end{proof}

\section{Simplicity and capstone (S2-E)}
\label{sec:sphere_simplicity}

\begin{lemma}

\leanok
[constant curvature branch: the centered circle realizes constant $\kappa$]
  \label{lem:spherical_circle_realizes}
  \lean{Gluck.sphericalCircle_realizes}
  \uses{def:realizes_spherical_curvature, lem:spherical_speed_circle}
  For every constant $c>0$ the centered circle $z(\theta)=-r^\ast\,i\,e^{i\theta}$ with
  $r^\ast=\sqrt{1+c^2}-c\in(0,1)$ is a simple closed curve realizing the constant
  curvature function $\kappa\equiv c$
  (\cref{def:realizes_spherical_curvature}). This discharges the constant disjunct of
  \cref{def:four_vertex_condition} in the capstone, exactly as the Euclidean
  \texttt{gluck\_converse} handles its constant branch with a round circle; the
  reconstruction flow and the degree argument are NOT needed for it.
\end{lemma}

\begin{proof}
  \uses{lem:spherical_speed_circle}
  \leanok
  $z$ is $C^1$ (an explicit exponential), $2\pi$-periodic, and
  $z'(\theta)=r^\ast e^{i\theta}\ne0$, so the tangent-angle witness is
  $\varphi(\theta)=\theta$ with $\|z'\|=r^\ast$. Confinement: $\|z\|=r^\ast<1$ since
  $\sqrt{1+c^2}<1+c$. The speed relation is the geodesic-circle identity: with
  $\langle z,ie^{i\theta}\rangle=-r^\ast$ the right-hand side is
  $(c+r^\ast)r^\ast$ and the left-hand side $\tfrac{1+r^{\ast2}}{2}$; they agree by the
  circle identity $1+r^{\ast2}=2r^\ast(c+r^\ast)$ (defining property of
  $r^\ast=\sqrt{1+c^2}-c$; this is the content of \cref{lem:spherical_speed_circle} read
  as an identity, with $c=(1-r^{\ast2})/(2r^\ast)$). Simplicity: on $[0,2\pi)$ the map
  $\theta\mapsto -r^\ast ie^{i\theta}$ is injective because $e^{i\theta}$ is.
\end{proof}

\begin{lemma}[spherical realization transfers under reparametrization]
  \label{lem:realizes_spherical_comp}
  \lean{Gluck.realizesSphericalCurvature_comp}
  \uses{def:realizes_spherical_curvature, lem:realizes_curvature_comp}
  If $z$ realizes the spherical curvature $\mu$
  (\cref{def:realizes_spherical_curvature}) and $\psi:\mathbb R\to\mathbb R$ is $C^1$
  with $\psi'>0$ everywhere, then $z\circ\psi$ realizes $\mu\circ\psi$.
\end{lemma}

\begin{proof}
  Mirror of the Euclidean \cref{lem:realizes_curvature_comp}
  (\texttt{Gluck.realizesCurvature\_comp}): the chain rule gives
  $(z\circ\psi)'=\psi'\cdot(z'\circ\psi)$, so the speed scales by $\psi'>0$, the
  tangent-angle witness is $\varphi\circ\psi$, and the spherical speed relation of
  \cref{def:realizes_spherical_curvature} at $\psi(t)$ multiplies through by
  $\psi'(t)>0$ on both sides — the extra conformal factor $(1+\|z\|^2)/2$ and the
  bracket term are pointwise in the curve value and transport unchanged. Confinement and
  regularity are pointwise. This is what pulls a realization of $\kappa\circ h_1$ back
  to a realization of $\kappa$ in the capstone: apply it with $\psi=h_1^{-1}$, whose
  existence, $C^1$ regularity, positive derivative, and circle-equivariance are exactly
  the PUBLIC in-tree lemma \texttt{Gluck.exists\_C1\_circle\_inverse}
  (\texttt{Reduction.lean}, \cref{lem:exists_c1_circle_inverse}) — reuse it directly,
  no local replication needed.
\end{proof}

\begin{lemma}[simplicity]
  \label{lem:spherical_simplicity}
  \lean{Gluck.spherical_simplicity}
  \uses{lem:spherical_endpoint_winding, lem:reconstruction_ode}
  The closing trajectory is simple: distinct parameters in $[0,2\pi)$ give distinct disk
  points, and (being confined to $\{|z|<1\}$) it is therefore simple on $S^2$.
\end{lemma}

\begin{proof}
  The disk representative is an ordinary plane curve with $z_\theta=q\,e^{i\theta}$, $q>0$,
  so the tangent direction $e^{i\theta}$ is strictly monotone in $\theta$; the disk metric
  plays no role in injectivity. The directional cosine (chord) integral controlling
  injectivity is therefore bounded away from $0$ exactly as in \cref{chap:simplicity}, via
  the public primitives \texttt{chord\_integral\_ne\_zero} /
  \texttt{isSimpleClosed\_reconstruct} (the Euclidean \texttt{simplicity\_transport} lemma
  is \texttt{private} in \texttt{DahlbergStep2} and is not importable, so simplicity is
  re-derived from the public chord primitive). Confinement to $\{|z|<1\}$ then transports
  simplicity to $S^2$ through the stereographic diffeomorphism. \emph{(Arc-level; S2-E.)}
\end{proof}
